%% Question bank chapter file - generated by bank_tools.py build (do not edit while the bank is still being built:
%% it is regenerated from the per-test data in Question Bank\data). Every question is one bankq block:
%%   \begin{bankq}[id=..., source={...}, paper=.., qno=.., type=..., subtopic={...}, answer={...}, answer_src=key|solved,
%%                 working={...}, difficulty=1-10, diff_parts={K S I C F}, diff_note={...},
%%                 group=<passage id or empty>, also_in={duplicate copies in other tests}]
%% difficulty = min(10, K+S+I+C+F-1): K concepts 1-3, S steps 1-4, I insight 0-2, C computation 0-2, F format 0-1
%%   ...question LaTeX, options, figures (\bfig / \bfigi -> img/)...
%%   \end{bankq}
%% Passages shared by several questions are bankpassage blocks; their questions carry group=<passage id>.
\documentclass[10pt]{article}
\usepackage[a4paper,margin=1.8cm]{geometry}
\usepackage[T1]{fontenc}
\usepackage{amsmath,amssymb,mathtools}
\usepackage{graphicx}
\usepackage[export]{adjustbox}
\usepackage{xcolor,needspace,keyval}
\usepackage{enumitem}
\usepackage{booktabs,array,tabularx}
\usepackage[hidelinks]{hyperref}
\setlength{\parindent}{0pt}
\setlength{\parskip}{3pt}
\makeatletter
\define@key{bq}{id}{\def\bq@id{#1}}\define@key{bq}{source}{\def\bq@src{#1}}\define@key{bq}{paper}{\def\bq@paper{#1}}
\define@key{bq}{qno}{\def\bq@qno{#1}}\define@key{bq}{type}{\def\bq@type{#1}}\define@key{bq}{subtopic}{\def\bq@sub{#1}}
\define@key{bq}{answer}{\def\bq@ans{#1}}\define@key{bq}{answer_src}{\def\bq@asrc{#1}}\define@key{bq}{working}{\def\bq@work{#1}}
\define@key{bq}{group}{\def\bq@group{#1}}\define@key{bq}{also_in}{\def\bq@also{#1}}
\define@key{bq}{difficulty}{\def\bq@diff{#1}}\define@key{bq}{diff_parts}{\def\bq@dparts{#1}}\define@key{bq}{diff_note}{\def\bq@dnote{#1}}
\newenvironment{bankq}[1][]{\def\bq@work{}\def\bq@also{}\def\bq@group{}\def\bq@diff{?}\def\bq@dparts{}\def\bq@dnote{}\setkeys{bq}{#1}%
  \par\needspace{4\baselineskip}\medskip\noindent{\small\textbf{\bq@id}\hfill\color{gray}[\bq@src, Paper \bq@paper, Q\bq@qno]
  \bq@type\ $\cdot$ \bq@sub\ $\cdot$ \textbf{D\bq@diff}}\par\nopagebreak\relax}
  {\par{\small\color{teal}Answer: \textbf{\bq@ans}\ifx\bq@work\@empty\else\ (solved: \bq@work)\fi
  \ifx\bq@also\@empty\else\par\color{gray}Also in: \bq@also\fi}\par}
\define@key{bp}{id}{\def\bp@id{#1}}\define@key{bp}{source}{\def\bp@src{#1}}\define@key{bp}{questions}{\def\bp@qs{#1}}
\newenvironment{bankpassage}[1][]{\setkeys{bp}{#1}\par\needspace{4\baselineskip}\bigskip\noindent
  {\small\textbf{Passage \bp@id}\hfill\color{gray}[\bp@src; questions \bp@qs]}\par\nopagebreak\relax\itshape}{\par}
\makeatother
\newcommand{\bfig}[2]{\par\smallskip{\centering\includegraphics[width=#1pt,max width=0.62\linewidth,max height=0.32\textheight]{img/#2}\par}\smallskip}
\newcommand{\bfigi}[2]{\includegraphics[width=#1pt,max width=\linewidth,max height=0.32\textheight,valign=c]{img/#2}}
\newcommand{\bmissing}[1]{\textit{[figure missing in the source conversion: \detokenize{#1}]}}
\newenvironment{qoptions}{\begin{itemize}[leftmargin=1.8em,itemsep=1pt,topsep=2pt,label={}]}{\end{itemize}}
\newcommand{\opt}[2]{\item[(#1)] #2}
\providecommand{\fwstart}{}\providecommand{\fwend}{}\providecommand{\qfullsplit}{}
\newcolumntype{L}[1]{>{\raggedright\arraybackslash}p{#1}}
\newcolumntype{C}[1]{>{\centering\arraybackslash}p{#1}}
\newcolumntype{R}[1]{>{\raggedleft\arraybackslash}p{#1}}
\providecommand{\sectionhdr}[1]{}\providecommand{\subsectionhdr}[1]{}


\begin{document}
{\Large\bfseries Area Under the Curve --- Question Bank}\par\smallskip
{\small Mathematics chapter M28. Questions from all converted test papers, grouped by source series. Answers: printed key where it matches the paper, otherwise solved (answer\_src=solved, with working). Exact repeats across papers appear once, with also\_in listing the other copies.}\par\medskip\hrule

\section*{Score-I papers}
\begin{bankq}[id=SI07-P1-MATH-S1-Q03, source={Score-I - Test 07 (04 Feb 2022)}, paper=1, qno=3, type=multi, subtopic={Area between curves}, answer={A,B}, answer_src=key, difficulty=5, diff_parts={K2 S2 I1 C1 F1 adj-1}, diff_note={concepts: area under x squared times power, one minus x | steps: S_k = beta function B(3, k+1); monotonic in k; sum; check statements}]
Let S$_k$ be the area bounded by the curve $y=x^2(1-x)^k$ and the lines $x=0$, $y=0$ and $x=1$. Then which of the following is/are correct.
\begin{qoptions}
\opt{A}{$\displaystyle\sum_{k=1}^{100}\mathrm S_k<\dfrac16$}
\opt{B}{If $\displaystyle\lim_{n\to\infty}\sum_{k=1}^{n}\mathrm S_k=\dfrac{p}{q}$, then $p+q=7$ where p and q are coprime numbers}
\opt{C}{$\mathrm S_{99}=\dfrac{1}{100\times101\times102}$}
\opt{D}{$\displaystyle\sum_{k=1}^{999}\mathrm S_k>\dfrac16$}
\end{qoptions}
\end{bankq}
\begin{bankq}[id=SI07-P1-MATH-S3-Q04, source={Score-I - Test 07 (04 Feb 2022)}, paper=1, qno=4, type=numerical, subtopic={Area between curves}, answer={0}, answer_src=key, difficulty=5, diff_parts={K2 S2 I1 C1 F1 adj-1}, diff_note={concepts: area enclosed by an inverse trigonometric curve, a line | steps: simplify the sum of two arcsin terms; curve becomes simple; area integral; compute}]
If area enclosed by curve $f(x)=\sin^{-1}\left(\dfrac{2|x|}{1+x^2}\right)+\sin^{-1}\left(\dfrac{|x|}{\sqrt{1+x^2}}\right)$ and $g(x)=\left|\tan^{-1}x\right|$ from x = $-1$ to x = 1 is $\pi-\ln b$. Then value of $[b-\pi]$ is, where [ ] is greatest integer function.
\end{bankq}
\begin{bankq}[id=SI09-P1-MATH-S1-Q02, source={Score-I - Test 09 (16 Feb 2022)}, paper=1, qno=2, type=single, subtopic={Area between curves}, answer={A}, answer_src=key, difficulty=4, diff_parts={K2 S2 I1 C1 F0 adj-1}, diff_note={concepts: area, a region with linear inequalities | steps: plot the region bounded by lines; vertices; compute area by trapezoid or integration}]
The area of the region\\
{\small$\left\{(x,y):\ 0\le x\le\dfrac94,\ 0\le y\le1,\ x\ge3y,\ x+y\ge2\right\}$}\\ is
\begin{qoptions}
\opt{A}{$\dfrac{11}{32}$}
\opt{B}{$\dfrac{35}{96}$}
\opt{C}{$\dfrac{37}{96}$}
\opt{D}{$\dfrac{13}{32}$}
\end{qoptions}
\end{bankq}
\begin{bankq}[id=SI11-P1-MATH-S1-Q03, source={Score-I - Test 11 (27 Feb 2022)}, paper=1, qno=3, type=single, subtopic={Area between curves}, answer={A}, answer_src=key, difficulty=3, diff_parts={K1 S2 I1 C1 F0 adj-1}, diff_note={concepts: area between two exponential curves in the first quadrant | steps: f for x > 0 is e^x - e^-x; intersection with g; integral of the difference; evaluate area}, also_in={JEE Advanced (2021 Series) - Test 02 (04 Feb 2021) P1 Q3}]
Let the functions $f:\mathbb R\to\mathbb R$ and $g:\mathbb R\to\mathbb R$ be defined by $f(x)=e^{x-1}-e^{-|x-1|}$ and $g(x)=\dfrac12\left(e^{x-1}+e^{1-x}\right)$. Then the area of the region in the first quadrant bounded by the curves $y=f(x)$, $y=g(x)$ and $x=0$ is
\begin{qoptions}
\opt{A}{$(2-\sqrt3)+\dfrac12(e-e^{-1})$}
\opt{B}{$(2+\sqrt3)+\dfrac12(e-e^{-1})$}
\opt{C}{$(2-\sqrt3)+\dfrac12(e+e^{-1})$}
\opt{D}{$(2+\sqrt3)+\dfrac12(e+e^{-1})$}
\end{qoptions}
\end{bankq}
\begin{bankq}[id=SI12-P1-MATH-S3-Q10, source={Score-I - Test 12 (09 Mar 2022)}, paper=1, qno=10, type=numerical, subtopic={Area between curves}, answer={4}, answer_src=key, difficulty=4, diff_parts={K1 S2 I1 C0 F1}, diff_note={concepts: area of a region defined with greatest integers | steps: integer pairs with product 2 including negatives; each is a unit square; total 4 | I: include the negative pairs}]
Find the area enclosed by the solution set of $[x]\cdot[y]=2$. Where $[\cdot]$ represent greatest integer function of $x$.
\end{bankq}
\begin{bankq}[id=SI18-P1-MATH-S2-Q07, source={Score-I - Test 18 (24 Apr 2022)}, paper=1, qno=7, type=multi, subtopic={Area between curves}, answer={A,C,D}, answer_src=key, difficulty=6, diff_parts={K2 S2 I2 C1 F1 adj-1}, diff_note={concepts: area between the logarithm curve axis, a horizontal line | steps: region from x = 0 to e - 1 under y = 1; area A = integral of (1 - ln(x+1)) or by inverse e^y - 1; A = e - 2; then floor and fractional part tests | I: integrate with respect to y and then evaluate greatest integer and fractional parts}]
If area enclosed by the graph of function $y=\ell n(x+1)$, $y$-axis and the line $y=1$ is $A$, then choose correct option(s)- (where [.] and \{.\} denotes greatest integer function and fractional part function respectively)
\begin{qoptions}
\opt{A}{\{A\} = e $-$ 2}
\opt{B}{[A] = 1}
\opt{C}{[A] = 0}
\opt{D}{[$-$A] = $-$1}
\end{qoptions}
\end{bankq}
\begin{bankq}[id=SI18-P2-MATH-S3-Q04, source={Score-I - Test 18 (24 Apr 2022)}, paper=2, qno=4, type=numerical, subtopic={Area between curves}, answer={4}, answer_src=key, difficulty=4, diff_parts={K2 S1 I1 C1 F1 adj-1}, diff_note={concepts: area between a function, its inverse | steps: curves symmetric about y = x; area = 2 times integral of (f(x) - x) from 0 to pi = 2 integral sin x = 4}, also_in={JEE Advanced (2021 Series) - Test 01 (24 Jan 2021) P2 Q4}]
Let $f:[0,\pi]\to[0,\pi]$ and $g:[0,\pi]\to[0,\pi]$ be two functions defined by $f(x)=x+\sin x$ and $g$ is inverse function of $f$. If $A$ is the area of the region bounded by $y=f(x)$ and $y=g(x)$, then the value of $A$ is
\end{bankq}
\begin{bankq}[id=SI26-P1-MATH-S1-Q02, source={Score-I - Test 26 (18 Dec 2022)}, paper=1, qno=2, type=multi, subtopic={Area between curves}, answer={A,D}, answer_src=key, difficulty=5, diff_parts={K1 S2 I1 C1 F1}, diff_note={concepts: bounds for the area under 2/(1+x^4) from 0 to 1 | steps: compare integrand with 2/(1+x^2) integrating to pi/2; compare with 2/(1+x^4) larger than the area of simple bounds; value ln 2 bound from another comparison; check four inequalities}]
Let us consider the function $f(x)=e^{-x}\cdot\sin x$, then
\begin{qoptions}
\opt{A}{Area bounded by the curve $f(x)$ and $x$-axis between ordinate $x=-2\pi$ and $x=0$ is $\left(\dfrac{e^{\pi}+1}{\sqrt2}\right)^2$}
\opt{B}{Area bounded by the curve $f(x)$ and $x$-axis between ordinate $x=-\pi$ and $x=\pi$ is $\dfrac{e^{\pi}-e^{-\pi}}{2}$}
\opt{C}{$\displaystyle\int_{m\pi}^{(m+1)\pi}f(x)\,dx<\int_{n\pi}^{(n+1)\pi}f(x)\,dx$ where $(m,n\in I)$ and $(m<n)$}
\opt{D}{$\displaystyle\int_{2m\pi}^{(2m+1)\pi}f(x)\,dx<\int_{2n\pi}^{(2n+1)\pi}f(x)\,dx$ where $(m,n\in I)$ and $(m>n)$}
\end{qoptions}
\end{bankq}
\begin{bankq}[id=SI26-P1-MATH-S1-Q05, source={Score-I - Test 26 (18 Dec 2022)}, paper=1, qno=5, type=multi, subtopic={Area between curves}, answer={C,D}, answer_src=key, difficulty=4, diff_parts={K1 S2 I1 C0 F1}, diff_note={concepts: least area between a parabola and a line with variable slope | steps: area between curves = (D)^(3/2)/6 with D from discriminant of x^2 - 3 - kx - 2; minimise over k; A least when k = 0; compute 3A/(10 sqrt5) | I: area of a parabolic segment through discriminant}]
If the least value of area bounded by $y=x^2-3$ and $y=kx+2$ be $A$, then the value of $\dfrac{3A}{10\sqrt5}$ is
\begin{qoptions}
\opt{A}{less than 2}
\opt{B}{greater than 5}
\opt{C}{less than 4}
\opt{D}{greater than 1}
\end{qoptions}
\end{bankq}
\begin{bankq}[id=SI26-P1-MATH-S4-Q07, source={Score-I - Test 26 (18 Dec 2022)}, paper=1, qno=7, type=numerical, subtopic={Area between curves}, answer={12}, answer_src=key, difficulty=7, diff_parts={K1 S3 I2 C1 F1}, diff_note={concepts: ratio of areas into which a floor function curve divides a circle | steps: circle centre (2,-1) radius 2; curve y = floor of a monotone expression takes integer steps; find segments cutting the circle; areas of circular segments and rectangles; ratio form gives k and m | I: step function boundaries inside the circle}]
The ratio of the areas in which the curve $y=\left[\dfrac{x^3}{101}+\dfrac{x}{36}\right]$ divides the circle $x^2+y^2-4x+2y+1=0$, (where $[.]$ denotes greatest integer function) is $\dfrac{k\pi-3\sqrt3}{m\pi+3\sqrt3}$ then value of $|k+m|$ \underline{\hspace{2.5em}}
\end{bankq}
\begin{bankq}[id=SI31-P2-MATH-S1-Q16, source={Score-I - Test 31 (29 Dec 2022)}, paper=2, qno=16, type=single, subtopic={Area between curves}, answer={B}, answer_src=key, difficulty=5, diff_parts={K2 S2 I1 C1 F0}, diff_note={concepts: area bounded by a circle, a polynomial of least degree from a limit condition | steps: limit (1 + f/x^3)^(1/x) = e requires f/x^4 tends to 1 so f = x^4 least degree with leading behaviour; find intersection with circle; area by integration}]
The smaller area bounded by the circle $x^2+y^2=2$ and $y=f(x)$ above X-axis where $f(x)$ is a polynomial of least degree satisfying $\displaystyle\lim_{x\to0}\left(1+\dfrac{f(x)}{x^3}\right)^{1/x}=e$ is
\begin{qoptions}
\opt{A}{$\dfrac\pi2-\dfrac35$}
\opt{B}{$\dfrac\pi2+\dfrac35$}
\opt{C}{$\dfrac\pi2+\dfrac85$}
\opt{D}{$\dfrac\pi2-\dfrac85$}
\end{qoptions}
\end{bankq}

\section*{Score-II papers}
\begin{bankq}[id=SII03-P1-MATH-S2-Q07, source={Score-II - Test 03 (05 May 2021)}, paper=1, qno=7, type=single, subtopic={Area between curves}, answer={1}, answer_src=key, difficulty=4, diff_parts={K1 S2 I1 C1 F0}, diff_note={concepts: least value of the area between two parabolas for parameter ranges | steps: intersection of the two parabolas depends on k1 and k2; area as an integral; monotonic in the parameters; evaluate the least value}]
Let `A' be the area bounded by the curves $y^2=4k_1x$, $k_1\in\left[\dfrac18,\dfrac14\right]$ and $x^2=4k_2(2-y)$, $k_2\in\left[\dfrac14,1\right]$ and y-axis in the first quadrant. Then the least value of A is
\end{bankq}
\begin{bankq}[id=SII05-P1-MATH-S1-Q14, source={Score-II - Test 05 (22 Mar 2022)}, paper=1, qno=14, type=single, subtopic={Area between curves}, answer={A}, answer_src=key, difficulty=4, diff_parts={K1 S2 I1 C1 F0}, diff_note={concepts: area common to a parabola region and a disc | steps: intersection of y^2 = 4x with x^2 + y^2 = 9/4; area of the parabola part plus circular segment; integrate; simplify to the given form}]
The area of the region\\ $\{(x,y):y^2\leq4x,\ 4x^2+4y^2\leq9\}$ is :-
\begin{qoptions}
\opt{A}{{\small $\dfrac{9\pi}8-\dfrac94\sin^{-1}\left(\dfrac13\right)+\dfrac{\sqrt2}6$}}
\opt{B}{{\small $\dfrac{9\pi}{16}-\dfrac98\sin^{-1}\left(\dfrac13\right)+\dfrac{\sqrt2}{12}$}}
\opt{C}{{\small $\dfrac{9\pi}4-\dfrac92\sin^{-1}\left(\dfrac13\right)+\dfrac{\sqrt2}3$}}
\opt{D}{{\small $\dfrac{9\pi}{32}-\dfrac9{16}\sin^{-1}\left(\dfrac13\right)+\dfrac{\sqrt2}{24}$}}
\end{qoptions}
\end{bankq}
\begin{bankq}[id=SII07-P1-MATH-S2-Q09, source={Score-II - Test 07 (10 Jul 2022)}, paper=1, qno=9, type=numerical, subtopic={Area between curves}, answer={Bonus}, answer_src=solved, working={Printed key shows letters for numerical-type questions here (misaligned) and the question was not solved; banked as Bonus.}, difficulty=5, diff_parts={K1 S2 I1 C1 F1}, diff_note={concepts: area enclosed by y = ln(2x + e) and the axes | steps: intercept x = (1-e)/2 negative; area by integrating x as a function of y; evaluate K and 1/K}]
If the area enclosed between the curve $y=\log_e(2x+e)$ and the coordinate axes is $K$, then $\dfrac1K$ is
\end{bankq}
\begin{bankq}[id=SII08-P2-MATH-S3-Q06, source={Score-II - Test 08 (05 Feb 2023)}, paper=2, qno=6, type=numerical, subtopic={Area between curves}, answer={7}, answer_src=key, difficulty=5, diff_parts={K1 S2 I1 C1 F1}, diff_note={concepts: area of a region bounded by piecewise functions | steps: intersect f with the triangular function g; region between curves for the range abs(x) at most 3/4; integrate; multiply by 9}]
Consider the functions $f,g:\mathbb{R}\to\mathbb{R}$ defined by $f(x) = x^2+\dfrac{5}{12}$ and
$$g(x) = \begin{cases}2\left(1-\dfrac{4|x|}{3}\right), & |x|\le\dfrac34,\\ 0, & |x|>\dfrac34.\end{cases}$$
If $\alpha$ is the area of the region\par\noindent
\resizebox{\linewidth}{!}{$\left\{(x,y)\in\mathbb{R}\times\mathbb{R}: |x|\le\dfrac34,\ 0\le y\le\min\{f(x),g(x)\}\right\}$}\par
then the value of $(3\alpha+5)$ is\_\_\_\_\_\_\_.
\end{bankq}
\begin{bankq}[id=SII12-P1-MATH-S2-Q07, source={Score-II - Test 12 (04 Feb 2024)}, paper=1, qno=7, type=multi, subtopic={Area between curves}, answer={A,B,C,D}, answer_src=key, difficulty=4, diff_parts={K1 S2 I1 C0 F1}, diff_note={concepts: areas above and below a curve and horizontal lines | steps: areas of the green and red regions split by a horizontal line of height h; use the intermediate value property of the area as a function of h; check each statement}]
Consider the square region $S=[0,1]\times[0,1]$.\\
Let $G=\{(x,y)\in S:y>f(x)\}$ be called the GREEN region\\
$R=\{(x,y)\in S:y<f(x)\}$ be called the RED region.\\
$L_h=\{(x,h)\in S:x\in[0,1]\}$ be the horizontal line drawn at a height $h\in[0,1]$.\\
If $f:[0,1]\to[0,1]$ be the function defined as
$$f(x)=\dfrac{x^3}{3}-x^2+\dfrac{5x}{9}+\dfrac{11}{36},$$
then which of the following statements is(are) true?
\begin{qoptions}
\opt{A}{There exists an $h\in\left[\dfrac16,\dfrac23\right]$ such that area of GREEN region above line $L_h$ equals that of GREEN region below line $L_h$.}
\opt{B}{There exists an $h\in\left[\dfrac16,\dfrac23\right]$ such that area of RED region above line $L_h$ equals that of RED region below line $L_h$.}
\opt{C}{There exists an $h\in\left[\dfrac16,\dfrac23\right]$ such that area of GREEN region above line $L_h$ equals that of RED region below line $L_h$.}
\opt{D}{There exists an $h\in\left[\dfrac16,\dfrac23\right]$ such that area of RED region above line $L_h$ equals that of GREEN region below line $L_h$.}
\end{qoptions}
\end{bankq}
\begin{bankq}[id=SII12-P1-MATH-S4-Q02, source={Score-II - Test 12 (04 Feb 2024)}, paper=1, qno=2, type=numerical, subtopic={Area between curves}, answer={6}, answer_src=key, difficulty=3, diff_parts={K1 S1 I1 C0 F1}, diff_note={concepts: area under a piecewise linear function and its maximum | steps: area = 4 from the geometry of triangles gives n; maximum value of f from the peaks}]
Let $n\in N$, $n>1$ and $f:[0,1]\to R$ be the function defined by
$$f(x)=\begin{cases}n(1-2nx), & 0\le x\le\dfrac{1}{2n},\\[4pt] 2n(2nx-1), & \dfrac{1}{2n}\le x\le\dfrac{3}{4n},\\[4pt] 4n(1-nx), & \dfrac{3}{4n}\le x\le\dfrac1n,\\[4pt] \dfrac{n}{n-1}(nx-1), & \dfrac1n\le x\le1.\end{cases}$$
If n is such that the area of the region bounded by the curves x = 0, x = 1, y = 0 and y = f(x) is 3 square units, then the maximum value of the function f is
\end{bankq}
\begin{bankq}[id=SII13-P1-MATH-S1-Q05, source={Score-II - Test 13 (06 Feb 2024)}, paper=1, qno=5, type=single, subtopic={Area between curves}, answer={A}, answer_src=key, difficulty=4, diff_parts={K1 S2 I1 C1 F0}, diff_note={concepts: area between a circle and a parabola | steps: intersection points of the circle and x^2 = 2y; region is inside the circle and outside the parabola; integrate in y; evaluate}]
Area of the region $\{(x,y):x^2+(y-2)^2\le4,\ x^2\ge2y\}$ is :
\begin{qoptions}
\opt{A}{$2\pi-\dfrac{16}{3}$}
\opt{B}{$\pi-\dfrac83$}
\opt{C}{$\pi+\dfrac83$}
\opt{D}{$2\pi+\dfrac{16}{3}$}
\end{qoptions}
\end{bankq}
\begin{bankq}[id=SII14-P1-MATH-S2-Q08, source={Score-II - Test 14 (06 Feb 2024)}, paper=1, qno=8, type=multi, subtopic={Area between curves}, answer={A,B}, answer_src=key, difficulty=5, diff_parts={K1 S2 I1 C1 F1}, diff_note={concepts: bounds for the area under 2/(1+x^4) from 0 to 1 | steps: compare integrand with 2/(1+x^2) integrating to pi/2; compare with 2/(1+x^4) larger than the area of simple bounds; value ln 2 bound from another comparison; check four inequalities}]
Let A be the area bounded by the curve $f(x)=\dfrac{2}{1+x^4}$, x-axis, y-axis and ordinate x = 1, then-
\begin{qoptions}
\opt{A}{$A>\dfrac\pi2$}
\opt{B}{$A>2\ell n2$}
\opt{C}{$0<A<\dfrac\pi2$}
\opt{D}{$\dfrac\pi2<A<2\ell n2$}
\end{qoptions}
\end{bankq}
\begin{bankq}[id=SII17-P1-MATH-S2-Q03, source={Score-II - Test 17 (13 Feb 2024)}, paper=1, qno=3, type=numerical, subtopic={Area between curves}, answer={8}, answer_src=solved, working={x=8/(y\^{}2+4): area=integral over y=8*(pi/2)=4pi; 8cos(2014*4pi)=8.}, difficulty=5, diff_parts={K1 S2 I1 C1 F1}, diff_note={concepts: area enclosed by a curve and evaluation of a cosine | steps: solve for x = 8/(y^2 + 4); area by integration with arctan; lambda; 8 cos(2014 lambda)}]
The area enclosed by the curve $xy^2=4(2-x)$ and y-axis is $\lambda$ square unit, then the value of $8\cos(2014\,\lambda)$ is equal to
\end{bankq}
\begin{bankq}[id=SII19-P1-MATH-S1-Q02, source={Score-II - Test 19 (20 Feb 2024)}, paper=1, qno=2, type=single, subtopic={Area between curves}, answer={C}, answer_src=key, difficulty=4, diff_parts={K1 S2 I1 C1 F0}, diff_note={concepts: area between a hyperbola and a line | steps: intersection points x y = 4; area by integrating; evaluate}]
The area enclosed by the curves $xy+4y=16$ and $x+y=6$ is equal to :
\begin{qoptions}
\opt{A}{$28-30\log_e2$}
\opt{B}{$30-28\log_e2$}
\opt{C}{$30-32\log_e2$}
\opt{D}{$32-30\log_e2$}
\end{qoptions}
\end{bankq}
\begin{bankq}[id=SII21-P2-MATH-S3-Q08, source={Score-II - Test 21 (14 Mar 2024)}, paper=2, qno=8, type=numerical, subtopic={Area between curves}, answer={7}, answer_src=key, difficulty=6, diff_parts={K2 S3 I1 C1 F1 adj-1}, diff_note={concepts: area under the maximum of a parabola, a tent function | steps: compare f and g: find the intersection points; area under f on the part where f is larger; area under g where g larger; integrate; 96 alpha and the unit digit}]
Consider the functions $f,g:\mathbb R\to\mathbb R$ defined by
\[f(x) = x^2+\dfrac{5}{12}\ \text{and}\ g(x) = \begin{cases}2\left(1-\dfrac{4|x|}{3}\right), & |x|\le\dfrac34,\\[4pt] 0, & |x|>\dfrac34.\end{cases}\]
If $\alpha$ is the area of the region
\[\left\{(x,y)\in\mathbb R\times\mathbb R : |x|\le\dfrac34,\ 0\le y\le\max\{f(x),g(x)\}\right\}\]
then the unit digit of $96\alpha$ is\underline{\hspace{2em}}.
\end{bankq}
\begin{bankq}[id=SII22-P1-MATH-S1-Q14, source={Score-II - Test 22 (14 Mar 2024)}, paper=1, qno=14, type=single, subtopic={Area between curves}, answer={B}, answer_src=key, difficulty=4, diff_parts={K1 S2 I1 C1 F0}, diff_note={concepts: area between a parabola and a line | steps: intersection of y^2 = 8x and y = sqrt2 x; area from x = 1 to the intersection; integral of the difference}]
The area of the region $S=\{(x,y):y^2\le8x,y\ge\sqrt2x,x\ge1\}$ is :
\begin{qoptions}
\opt{A}{$\dfrac{13\sqrt2}{6}$}
\opt{B}{$\dfrac{11\sqrt2}{6}$}
\opt{C}{$\dfrac{5\sqrt2}{6}$}
\opt{D}{$\dfrac{19\sqrt2}{6}$}
\end{qoptions}
\end{bankq}
\begin{bankq}[id=SII24-P1-MATH-S1-Q18, source={Score-II - Test 24 (17 Mar 2024)}, paper=1, qno=18, type=single, subtopic={Area between curves}, answer={C}, answer_src=key, difficulty=4, diff_parts={K1 S2 I1 C1 F0}, diff_note={concepts: smaller area between curves with square roots | steps: curve sqrt x + sqrt abs y = 1 and the diamond abs x + abs y = 1; region between them in the first quadrant; integrate; double}]
The smaller area (in sq. units) included between the curves $\sqrt x+\sqrt{|y|}=1$ and $|x|+|y|=1$ is
\begin{qoptions}
\opt{A}{$\dfrac13$}
\opt{B}{$\dfrac43$}
\opt{C}{$\dfrac23$}
\opt{D}{$\dfrac53$}
\end{qoptions}
\end{bankq}
\begin{bankq}[id=SII24-P1-MATH-S1-Q19, source={Score-II - Test 24 (17 Mar 2024)}, paper=1, qno=19, type=single, subtopic={Area between curves}, answer={A}, answer_src=key, difficulty=4, diff_parts={K1 S2 I1 C1 F0}, diff_note={concepts: area between a sine curve and a line through points | steps: line passes through (pi/6, 1/2) and (5pi/6, 1/2)?; intersection points; area by integrating sin x minus line}]
Area bounded by $y=\sin x$ and $y-\dfrac12=\dfrac3{2\pi}\left(x-\dfrac\pi6\right)$ is equal to ($x>0$)
\begin{qoptions}
\opt{A}{$\dfrac{\sqrt3}2-\dfrac\pi4$}
\opt{B}{$\sqrt3-\dfrac\pi2$}
\opt{C}{$\dfrac{\sqrt3}2-\dfrac\pi6$}
\opt{D}{$2\sqrt3-\pi$}
\end{qoptions}
\end{bankq}
\begin{bankq}[id=SII25-P1-MATH-S1-Q14, source={Score-II - Test 25 (20 Mar 2024)}, paper=1, qno=14, type=single, subtopic={Area between curves}, answer={A}, answer_src=key, difficulty=4, diff_parts={K1 S2 I1 C1 F0}, diff_note={concepts: area bounded by the constant solution and the logarithm | steps: y = 2 and y = ln x meet at x = e^2; area from y axis: integral of (2 - ln x) from; compute with the correct region}]
The area bounded by the line $x=1$ and the curve $\sqrt{\dfrac yx}+\sqrt{\dfrac xy}=4$ is :
\begin{qoptions}
\opt{A}{$4\sqrt3$}
\opt{B}{$\sqrt3$}
\opt{C}{$2\sqrt3$}
\opt{D}{$3\sqrt3$}
\end{qoptions}
\end{bankq}
\begin{bankq}[id=SII25-P2-MATH-S1-Q05, source={Score-II - Test 25 (20 Mar 2024)}, paper=2, qno=5, type=multi, subtopic={Area between curves}, answer={A,C}, answer_src=key, difficulty=7, diff_parts={K1 S3 I2 C1 F1}, diff_note={concepts: areas of nested diamonds with recursive centres and sizes | steps: abs(x-a) + abs(y) = b is a square of diagonal 2b with area 2 b^2; sequences a_i b_i; areas form a geometric series 2 b_i^2; total area and overlap of consecutive diamonds; check statements | I: area of the diamond is 2 b^2}]
If $A_i$ represents the area bounded by $|x-a_i|+|y|=b_i$ ; $i\in$ natural number and $a_{i+1}=a_i+\dfrac32b_i$ and $b_{i+1}=\dfrac{b_i}2$ ; $a_1=0$ ; $b_1=32$ then
\begin{qoptions}
\opt{A}{$A_3=128$}
\opt{B}{$\displaystyle\lim_{n\to\infty}\sum_{i=1}^{n}A_i=\dfrac43(16)^2$}
\opt{C}{$\displaystyle\lim_{n\to\infty}\sum_{i=1}^{n}A_i=\dfrac83(32)^2$}
\opt{D}{$A_3=256$}
\end{qoptions}
\end{bankq}
\begin{bankq}[id=SII27-P1-MATH-S1-Q11, source={Score-II - Test 27 (23 Mar 2024)}, paper=1, qno=11, type=single, subtopic={Area between curves}, answer={B}, answer_src=key, difficulty=4, diff_parts={K1 S2 I1 C1 F0}, diff_note={concepts: area bounded by a curve  its inverse region | steps: region bounded by sqrt(1+x^2) the line y=x y axis and y=-x+a; inverse function symmetric about y=x; area relations between the original and inverse regions; area of the inverse region in terms of k}]
The area of the region between the curves $y=\sqrt{\dfrac{1+\sin x}{\cos x}}$ and $y=\sqrt{\dfrac{1-\sin x}{\cos x}}$ bounded by the lines $x=0$ and $x=\dfrac\pi4$ is
\begin{qoptions}
\opt{A}{$\displaystyle\int_{\sqrt2+1}^{\sqrt2-1}\dfrac{t}{(1+t^2)\sqrt{1-t^2}}\,dt$}
\opt{B}{$\displaystyle\int_0^{\sqrt2-1}\dfrac{4t}{(1+t^2)\sqrt{1-t^2}}\,dt$}
\opt{C}{$\displaystyle\int_0^{\sqrt2+1}\dfrac{4t}{(1+t^2)\sqrt{1-t^2}}\,dt$}
\opt{D}{$\displaystyle\int_0^{\sqrt2+1}\dfrac{t}{(1+t^2)\sqrt{1-t^2}}\,dt$}
\end{qoptions}
\end{bankq}
\begin{bankq}[id=SII30-P1-MATH-S2-Q03, source={Score-II - Test 30 (27 Mar 2024)}, paper=1, qno=3, type=numerical, subtopic={Area between curves}, answer={Bonus}, answer_src=solved, working={Printed key shows letters for numeric-type items here (misaligned) and the question was not solved; banked as Bonus.}, difficulty=6, diff_parts={K1 S3 I1 C1 F1}, diff_note={concepts: area of region bounded by line and parabolas | steps: sketch the region; split the area by x ranges; integrate; reduce m by n; m + n}]
Let the area of the region $\{(x,y):x-2y+4\geq0,\ x+2y^2\geq0,\ x+4y^2\leq8,\ y\geq0\}$ be $\dfrac mn$, where $m$ and $n$ are coprime numbers. Then $m+n$ is equal to \underline{\hspace{2em}}.
\end{bankq}
\begin{bankq}[id=SII30-P1-MATH-S2-Q05, source={Score-II - Test 30 (27 Mar 2024)}, paper=1, qno=5, type=numerical, subtopic={Area between curves}, answer={Bonus}, answer_src=solved, working={Printed key shows letters for numeric-type items here (misaligned) and the question was not solved; banked as Bonus.}, difficulty=5, diff_parts={K1 S2 I1 C1 F1}, diff_note={concepts: area under minimum of a parabola and a line | steps: find where x^2 + 2 = 2x + 2 at 0 and 2; area below the lower curve in two parts; 12 A}]
Let the area of the region $\{(x,y):0\leq x\leq3,\ 0\leq y\leq\min\{x^2+2,\,2x+2\}\}$ be A. Then 12A is equal to \underline{\hspace{2em}}.
\end{bankq}
\begin{bankpassage}[id=SII36-P1-MATH-S3-G5-6, source={Score-II - Test 36 (07 Feb 2025)}, questions={5-6}]
If the area bounded by the function $f(x)=\cos^{-1}(\cos x)^2$ and $g(x)=\left(\cos^{-1}(\cos(x))\right)^2$ and $x=\pm2\pi$ is equal to $\dfrac{a(\pi^3+1)-b\pi^2}c$, where $a,b,c\in\mathbb N$, where a \& c are co-prime.
\end{bankpassage}
\begin{bankq}[id=SII36-P1-MATH-S3-Q05, source={Score-II - Test 36 (07 Feb 2025)}, paper=1, qno=5, type=numerical, subtopic={Area between curves}, answer={0.5}, answer_src=key, difficulty=6, diff_parts={K1 S3 I1 C1 F1}, diff_note={concepts: area between a function of inverse cosine of cos and its square | steps: sketch (arccos(cos x))^2 and arccos(cos^2 x); points of intersection; area on [0, 2 pi] by symmetry doubled; integrate the quadratic and linear pieces; compare with the stated form to read a b c}, group=SII36-P1-MATH-S3-G5-6]
The value of $\dfrac1{12}\left(a+\dfrac bc\right)$ is
\end{bankq}
\begin{bankq}[id=SII36-P1-MATH-S3-Q06, source={Score-II - Test 36 (07 Feb 2025)}, paper=1, qno=6, type=numerical, subtopic={Area between curves}, answer={3.5}, answer_src=key, difficulty=5, diff_parts={K2 S2 I0 C1 F1}, diff_note={concepts: probability of a sum divisible by a number, area result | steps: get a b and c from the area computation; value a+b+c; count the sums of three dice divisible by it; probability times 36}, group=SII36-P1-MATH-S3-G5-6]
Three dice are thrown simultaneously and p is the probability that sum of numbers appearing on each of the three dice is divisible by $a+b+c$. Then $36p$ is equal to
\end{bankq}
\begin{bankq}[id=SII39-P2-MATH-S3-Q01, source={Score-II - Test 39 (02 Mar 2025)}, paper=2, qno=1, type=numerical, subtopic={Area between curves}, answer={4}, answer_src=key, difficulty=6, diff_parts={K1 S3 I1 C1 F1}, diff_note={concepts: area of the loop of 3y^2 = x^2(3-x) | steps: loop for x from 0 to 3; y = x sqrt(3 - x)/sqrt3; area = 2 integral; substitute u = 3 - x; evaluate; floor}]
The area of the loop of the curve $3y^2=x^2(3-x)$ is `A' then find [A] where [ ] represent G.I.F.
\end{bankq}
\begin{bankq}[id=SII40-P1-MATH-S2-Q07, source={Score-II - Test 40 (16 Mar 2025)}, paper=1, qno=7, type=multi, subtopic={Area between curves}, answer={A,B,C}, answer_src=key, difficulty=4, diff_parts={K1 S2 I1 C0 F1}, diff_note={concepts: areas above and below a curve and horizontal lines | steps: areas of the green and red regions split by a horizontal line of height h; use the intermediate value property of the area as a function of h; check each statement}]
Let $f:[0,1]\to[0,1]$ be the function defined by $36f(x)=12x^3-36x^2+20x+17$. Consider the square region $S=[0,1]\times[0,1]$. Let $G=\{(x,y)\in S:y>f(x)\}$ be called the green region and $R=\{(x,y)\in S:y<f(x)\}$ be called the red region. Let $L_h=\{(x,h)\in S:x\in[0,1]\}$ be the horizontal line drawn at a height $h\in[0,1]$. Then which of the following statements is(are) true ?
\begin{qoptions}
\opt{A}{There exists an $h\in\left[\dfrac14,\dfrac23\right]$ such that the area of the red region above the line $L_h$ equals the area of the green region below the line $L_h$}
\opt{B}{There exists an $h\in\left[\dfrac14,\dfrac23\right]$ such that the area of the red region above the line $L_h$ equals the area of the red region below the line $L_h$}
\opt{C}{There exists an $h\in\left[\dfrac14,\dfrac23\right]$ such that the area of the green region above the line $L_h$ equals the area of the red region below the line $L_h$}
\opt{D}{There exists an $h\in\left[\dfrac14,\dfrac23\right]$ such that the area of the green region above the line $L_h$ equals the area of the green region below the line $L_h$}
\end{qoptions}
\end{bankq}
\begin{bankq}[id=SII40-P1-MATH-S4-Q02, source={Score-II - Test 40 (16 Mar 2025)}, paper=1, qno=2, type=numerical, subtopic={Area between curves}, answer={12}, answer_src=key, difficulty=3, diff_parts={K1 S1 I1 C0 F1}, diff_note={concepts: area under a piecewise linear function and its maximum | steps: area = 4 from the geometry of triangles gives n; maximum value of f from the peaks}]
Let $n\geq2$ be a natural number and $f:[0,1]\to\mathbb R$ be the function defined by\\
$f(x)=\begin{cases}n(1-2nx)&\text{if }0\leqslant x\leqslant\frac1{2n}\\ 2n(2nx-1)&\text{if }\frac1{2n}\leqslant x\leqslant\frac3{4n}\\ 4n(1-nx)&\text{if }\frac3{4n}\leqslant x\leqslant\frac1n\\ \frac n{n-1}(nx-1)&\text{if }\frac1n\leqslant x\leqslant1\end{cases}$\\
If n is such that the area of the region bounded by the curves x = 0 , x = 1 , y = 0 and y = f (x) is 6, then the maximum value of the function f is
\end{bankq}
\begin{bankq}[id=SII43-P1-MATH-S1-Q19, source={Score-II - Test 43 (23 Mar 2025)}, paper=1, qno=19, type=single, subtopic={Area between curves}, answer={C}, answer_src=key, difficulty=5, diff_parts={K1 S3 I1 C1 F0}, diff_note={concepts: area between a solution of a differential equation  a piecewise max function | steps: f(x) = 2 (solution y = 2); g(x) is the max of three functions; plot pieces; area between y = 2 and g on the y axis boundary to their intersection; integrate pieces}]
The area enclosed by $y=\ell nx$, $y=\ell n|x|$, $y=|\ell nx|$ and $y=|\ell n|x||$ is
\begin{qoptions}
\opt{A}{5}
\opt{B}{2}
\opt{C}{4}
\opt{D}{1}
\end{qoptions}
\end{bankq}
\begin{bankq}[id=SII44-P1-MATH-S1-Q04, source={Score-II - Test 44 (23 Mar 2025)}, paper=1, qno=4, type=single, subtopic={Area between curves}, answer={Bonus}, answer_src=solved, working={Printed key covers only one subject (chemistry) and this question was not solved; banked as Bonus.}, difficulty=4, diff_parts={K1 S2 I2 C0 F0}, diff_note={concepts: area under a double exponential curve by recognising a derivative | steps: the integrand is the derivative of 2^(2^(x+...)) times constants; the area from 1 to m telescopes to a difference of two powers of 2; A(m)+2 is evaluated at several m and classified as prime or composite | I: the integrand is the exact derivative of a double exponential so the area telescopes}]
Let A(m) denotes the area of the region bounded by the curve $y=2^{2^{x+x-2}}\cdot(\ln2)^2$ and x-axis between the lines x = 1 and x = m (m$\in\mathbb N$, m $>$ 2), then A(m) + 2 is

\begin{qoptions}
\opt{A}{A prime number for all m}
\opt{B}{A composite number for all m}
\opt{C}{A prime number for some m}
\opt{D}{An irrational number for some m}
\end{qoptions}
\end{bankq}
\begin{bankq}[id=SII45-P1-MATH-S1-Q18, source={Score-II - Test 45 (26 Mar 2025)}, paper=1, qno=18, type=single, subtopic={Area between curves}, answer={A}, answer_src=key, difficulty=5, diff_parts={K1 S3 I1 C1 F0}, diff_note={concepts: area of the region x^10 - x^4 + y^2 <= 0 | steps: y^2 <= x^4 - x^10; area = 2 int of sqrt; substitution; numeric; greatest integer}]
Area of the region in $x$-$y$ plane satisfying $x^{10}-x^4+y^2\leq0$, is $A$, then $[A]$ (where $[x]$ represents greatest integer function)
\begin{qoptions}
\opt{A}{1}
\opt{B}{2}
\opt{C}{3}
\opt{D}{4}
\end{qoptions}
\end{bankq}
\begin{bankq}[id=SII47-P1-MATH-S1-Q11, source={Score-II - Test 47 (30 Mar 2025)}, paper=1, qno=11, type=single, subtopic={Area between curves}, answer={D}, answer_src=key, difficulty=2, diff_parts={K1 S1 I1 C0 F0}, diff_note={concepts: area of a region defined by modulus conditions | steps: region between two squares abs(x-2)+abs(y+1) between 1 and 2; area = 8 - 2 = 6}]
The area of the region A represented by $A=\{(x,y)\mid(x,y)\in R\times R,\ 1\le|x-2|+|y+1|\le2\}$ is
\begin{qoptions}
\opt{A}{9 sq. units}
\opt{B}{4 sq. units}
\opt{C}{5 sq. units}
\opt{D}{6 sq. units}
\end{qoptions}
\end{bankq}

\section*{Score Advanced papers}
\begin{bankq}[id=SA02-P2-MATH-S1-Q08, source={Score Advanced - Test 02 (13 May 2021)}, paper=2, qno=8, type=multi, subtopic={Area between curves}, answer={B,C,D}, answer_src=solved, working={Region x in (-1,1/2), -1-x<=y<0: area 9/8}, difficulty=6, diff_parts={K2 S2 I2 C1 F1 adj-1}, diff_note={concepts: area of a region with radicals, area between curves | steps: domain conditions; region; area}, also_in={JEE Advanced (2021 Series) - Test 10 (13 May 2021) P2 Q8}]
Area bounded by $\sqrt{x + y + 1} < \sqrt{y - x + 2}$ \& y $<$ 0 is less than
\begin{qoptions}
\opt{A}{$\dfrac{9}{16}$}
\opt{B}{$\dfrac{9}{5}$}
\opt{C}{$\dfrac{5}{4}$}
\opt{D}{$\dfrac{16}{5}$}
\end{qoptions}
\end{bankq}
\begin{bankq}[id=SA03-P1-MATH-S1-Q03, source={Score Advanced - Test 03 (27 Jun 2021)}, paper=1, qno=3, type=multi, subtopic={Area between curves}, answer={A,D}, answer_src=solved, working={Rhombus with diagonals 6 and 4: area 12}, difficulty=4, diff_parts={K1 S2 I1 C0 F1}, diff_note={concepts: area bounded by a modulus curve | steps: rhombus with the given intercepts; area; compare with options}]
Area bounded by the curve 2$|$x$|$ + 3$|$y$|$ $\leq$ 6 is n, then
\begin{qoptions}
\opt{A}{n $\geq$ 12}
\opt{B}{n $\leq$ 10}
\opt{C}{n $>$ 12}
\opt{D}{n $>$ 10}
\end{qoptions}
\end{bankq}
\begin{bankq}[id=SA03-P1-MATH-S1-Q06, source={Score Advanced - Test 03 (27 Jun 2021)}, paper=1, qno=6, type=multi, subtopic={Area between curves}, answer={A,B,C}, answer_src=solved, working={S=(e\^{}(a/2)-1)\^{}2 so S/a\^{}2 tends to 1/4 and n=4}, difficulty=6, diff_parts={K1 S2 I2 C1 F1}, diff_note={concepts: area between exponential curves | steps: intersection point; area as a function of a; maximum or limit; check}]
Consider two curves y = e$^x$, y = e$^{a-x}$, where a is positive number. Let S be the area of the part surrounded by C$_1$,C$_2$ and the y-axis, if $\displaystyle\lim_{a \to 0}\dfrac{S}{a^2}$ is $\dfrac{1}{n}$, then
\begin{qoptions}
\opt{A}{n is the first even composite number}
\opt{B}{number of positive divisors of n is 3}
\opt{C}{n is a perfect square}
\opt{D}{n is an odd number}
\end{qoptions}
\end{bankq}
\begin{bankq}[id=SA03-P2-MATH-S2-Q04, source={Score Advanced - Test 03 (27 Jun 2021)}, paper=2, qno=4, type=numerical, subtopic={Area between curves}, answer={5.33}, answer_src=solved, working={16/3}, difficulty=4, diff_parts={K1 S2 I1 C0 F1}, diff_note={concepts: area between two parabolas | steps: intersection points; integrate difference; compute}]
The area bounded by $y^2=4x$ and $x^2=4y$ is
\end{bankq}
\begin{bankq}[id=SA03-P2-MATH-S2-Q05, source={Score Advanced - Test 03 (27 Jun 2021)}, paper=2, qno=5, type=numerical, subtopic={Area between curves}, answer={50}, answer_src=solved, working={7/120 so b-10a=50}, difficulty=5, diff_parts={K1 S2 I1 C1 F1}, diff_note={concepts: area under a quartic between minima | steps: find minima; integrate absolute value; compute fraction}]
The area between the curve $y=2x^4-x^2$, the x-axis and the ordinates of two minima of the curve is of the form $\dfrac{a}{b}$ (where a and b are co-prime). Then the value of $(b-10a)$ is-
\end{bankq}
\begin{bankq}[id=SA03-P2-MATH-S2-Q07, source={Score Advanced - Test 03 (27 Jun 2021)}, paper=2, qno=7, type=numerical, subtopic={Area between curves}, answer={2.75}, answer_src=solved, working={(19 sqrt3-30)/12 so (3+30)/12=2.75}, difficulty=6, diff_parts={K1 S2 I2 C1 F1}, diff_note={concepts: area between two curves with trigonometric functions | steps: intersection; integrate difference using identities; evaluate | I: symmetry about pi/4}]
Area (in square units) bounded by the curves $f(x)=(\cos x+\sec x)^2$ and $g(x)=(\sin x+\mathrm{cosec}\,x)^2$ between $x=\dfrac{\pi}{4}$ and $x=\dfrac{\pi}{3}$ is of the form $\dfrac{19\sqrt{b}-c}{d}$, s.t. H.C.F of b,c and d is 3, then the value of $\dfrac{b+c}{d}$ is-
\end{bankq}
\begin{bankq}[id=SA03-P2-MATH-S2-Q08, source={Score Advanced - Test 03 (27 Jun 2021)}, paper=2, qno=8, type=numerical, subtopic={Area between curves}, answer={2.5}, answer_src=solved, working={4.5-2}, difficulty=5, diff_parts={K1 S2 I1 C1 F1}, diff_note={concepts: area bounded by pair of lines and axes | steps: separate the lines; vertices; area of the triangle}]
The area bounded by pair of straight lines $x^2+y^2-2xy-5x+5y+6=0$ and the co-ordinate axes is-
\end{bankq}
\begin{bankq}[id=SA05-P1-MATH-S1-Q16, source={Score Advanced - Test 05 (15 Jul 2021)}, paper=1, qno=16, type=single, subtopic={Area between curves}, answer={D}, answer_src=key, difficulty=3, diff_parts={K1 S1 I1 C1 F0}, diff_note={concepts: area of a circular segment | steps: region outside the line inside the unit circle}]
Area bounded by x$^2$ + y$^2$ $\leq$ 1 $\leq$ x + y is equal to-
\begin{qoptions}
\opt{A}{$\dfrac{\pi^2}{4}$}
\opt{B}{$\dfrac{\pi^2}{2}$}
\opt{C}{$\dfrac{\pi}{4}+\dfrac{1}{2}$}
\opt{D}{$\dfrac{\pi}{4}-\dfrac{1}{2}$}
\end{qoptions}
\end{bankq}
\begin{bankq}[id=SA10-P2-MATH-S3-Q04, source={Score Advanced - Test 10 (06 May 2023)}, paper=2, qno=4, type=numerical, subtopic={Area between curves}, answer={4}, answer_src=key, difficulty=8, diff_parts={K1 S4 I2 C1 F1}, diff_note={concepts: area between a curve and its inverse | steps: reflect about y = x; symmetric area; integrate; compute; reflect about y = x; intersection points; area under the curve; area between the curve and the line; double the area; evaluate; compute | I: area by symmetry about y=x}]
The area bounded by the curve f(x) = x + sinx and its inverse between the ordinates x = 0 to x = 3$\pi$ is 3K, K $\in$ N then K is
\end{bankq}
\begin{bankq}[id=SA11-P1-MATH-S3-Q05, source={Score Advanced - Test 11 (09 May 2023)}, paper=1, qno=5, type=numerical, subtopic={Area between curves}, answer={8}, answer_src=key, difficulty=5, diff_parts={K1 S2 I1 C1 F1}, diff_note={concepts: region defined by quadratic inequality | steps: complete the square; shape of S; area; compute}, also_in={Score Advanced - Test 29 (01 May 2024) P1 Q5}]
Let $f$ (x) = x$^2$ + 6x + 1 and S denotes the set of points (x, y) in the coordinate plane such that $f$(x) + $f$(y) $\leq$ 0 and $f$(x) $-$ $f$(y) $\leq$ 0. The area of S is k$\pi$ then k =
\end{bankq}
\begin{bankq}[id=SA14-P1-MATH-S2-Q09, source={Score Advanced - Test 14 (16 May 2023)}, paper=1, qno=9, type=single, subtopic={Area between curves}, answer={C}, answer_src=key, difficulty=4, diff_parts={K1 S2 I1 C1 F0}, diff_note={concepts: areas bounded by curves matched  values | steps: find intersection points; integrate piecewise; areas of three regions; match}]
Match the following:\par
{\footnotesize\setlength{\tabcolsep}{1.5pt}\noindent\begin{tabular}{|c|L{4.8cm}|c|c|}
\hline
 & List-I & & List-II \\
\hline
(P) & Area bounded by y = x$^2$ + 2 and y = 2 $|$x$|$ $-$ cos$\pi$x is equal to A, then 3A equals & (1) & 4 \\
\hline
(Q) & The value of 2cot(cot$^{-1}$3 + cot$^{-1}$7 + cot$^{-1}$13 + cot$^{-1}$21) & (2) & $\dfrac{8}{3}$ \\
\hline
(R) & Tangent and normal at the ends A \& C of focal chord AC of parabola y$^2$ = 4x intersect at B and D, then minimum area of ABCD is & (3) & 3 \\
\hline
(S) & Number of integral values of ordered pair ($\alpha$, $\beta$) for which the area common to x$^2$ + y$^2$ $-$ 2$\alpha^2$x $-$ 2$\beta^2$y + c = 0 and its image in x + y = 1 is maximum is & (4) & 8 \\
\hline
\end{tabular}}\par
Code:
\begin{qoptions}
\opt{A}{P $\rightarrow$ 4;Q $\rightarrow$ 3;R $\rightarrow$ 2;S $\rightarrow$ 3}
\opt{B}{P $\rightarrow$ 1;Q $\rightarrow$ 2;R $\rightarrow$ 1;S $\rightarrow$ 3}
\opt{C}{P $\rightarrow$ 4;Q $\rightarrow$ 3;R $\rightarrow$ 4;S $\rightarrow$ 1}
\opt{D}{P $\rightarrow$ 1;Q $\rightarrow$ 3;R $\rightarrow$ 4;S $\rightarrow$ 2}
\end{qoptions}
\end{bankq}
\begin{bankq}[id=SA16-P1-MATH-S2-Q08, source={Score Advanced - Test 16 (20 May 2023)}, paper=1, qno=8, type=multi, subtopic={Area between curves}, answer={C,D}, answer_src=key, difficulty=7, diff_parts={K1 S3 I1 C2 F1}, diff_note={concepts: region x^2 <= y^2 <= 2 - x: extremes  lattice points | steps: region bounded by the lines y = x and y = -x and the parabola y^2 = 2 - x; maximum of x^2 + y^2 at vertices; count integral points; area; check statements | C: region definition by squared inequalities with lattice counting}]
Consider the region $R:x^2\leq y^2\leq2-x$, then which of the following is/are \textbf{FALSE} ?
\begin{qoptions}
\opt{A}{Maximum value of $x^2+y^2$ is 8, $\forall(x,y)\in R$}
\opt{B}{Number of integral points (both coordinates integer) in $R$ is 9}
\opt{C}{Area of $R$ is $\dfrac{13}{6}$ sq.units}
\opt{D}{Minimum value of $x^2+y^2$ is 4, $\forall(x,y)\in R$}
\end{qoptions}
\end{bankq}
\begin{bankpassage}[id=SA16-P2-MATH-S2-G7-8, source={Score Advanced - Test 16 (20 May 2023)}, questions={7-8}]
Let $y=f(x)$ be a solution of the differential equation $\dfrac{dy}{dx}=-4x(y-2)$ with $f(0)=2$ and $g(x)=\max.\left\{1+\dfrac x2-\left[\dfrac x2\right],\sin x,|x-1|\right\}$, $h(x)=\ln x$\par
[\textbf{Note:} $[y]$ denotes greatest integer function less than or equal to $y$]\par
\end{bankpassage}
\begin{bankq}[id=SA16-P2-MATH-S2-Q07, source={Score Advanced - Test 16 (20 May 2023)}, paper=2, qno=7, type=single, subtopic={Area between curves}, answer={A}, answer_src=key, difficulty=5, diff_parts={K1 S3 I1 C1 F0}, diff_note={concepts: area between a solution of a differential equation  a piecewise max function | steps: f(x) = 2 (solution y = 2); g(x) is the max of three functions; plot pieces; area between y = 2 and g on the y axis boundary to their intersection; integrate pieces}, group=SA16-P2-MATH-S2-G7-8]
Area enclosed by the curves $y=f(x)$, $y=g(x)$ and the y-axis, is
\begin{qoptions}
\opt{A}{1}
\opt{B}{2}
\opt{C}{3}
\opt{D}{4}
\end{qoptions}
\end{bankq}
\begin{bankq}[id=SA16-P2-MATH-S2-Q08, source={Score Advanced - Test 16 (20 May 2023)}, paper=2, qno=8, type=single, subtopic={Area between curves}, answer={B}, answer_src=key, difficulty=4, diff_parts={K1 S2 I1 C1 F0}, diff_note={concepts: area bounded by the constant solution and the logarithm | steps: y = 2 and y = ln x meet at x = e^2; area from y axis: integral of (2 - ln x) from; compute with the correct region}, group=SA16-P2-MATH-S2-G7-8]
Area bounded by the curve $y=f(x)$, $y=h(x)$ and the y-axis, is
\begin{qoptions}
\opt{A}{$e$}
\opt{B}{$e^2$}
\opt{C}{$\dfrac{3e}{2}$}
\opt{D}{$e^2-1$}
\end{qoptions}
\end{bankq}
\begin{bankq}[id=SA19-P1-MATH-S2-Q09, source={Score Advanced - Test 19 (04 Apr 2024)}, paper=1, qno=9, type=multi, subtopic={Area between curves}, answer={A,B,D}, answer_src=key, difficulty=6, diff_parts={K1 S2 I2 C1 F1}, diff_note={concepts: area under y = ln^4 x  1 to e^4 | steps: integrate by parts repeatedly: integral of ln^4 x dx = x(ln^4 - 4 ln^3 + 12 ln^2 - 24 ln + 24); evaluate between 1 and e^4; area; compare with bounds | I: reduction formula for powers of the logarithm}]
Let A denote the area bounded by $y=\ell n^4x$ and lines $y=0$ and $x=e^4$, then -
\begin{qoptions}
\opt{A}{$A<256e^4-175e^3-81$}
\opt{B}{$A<256e^4-256$}
\opt{C}{$A>256\left(e^4-1\right)$}
\opt{D}{$A>\dfrac{337}{2}e^3(e-1)$}
\end{qoptions}
\end{bankq}
\begin{bankq}[id=SA20-P1-MATH-S2-Q09, source={Score Advanced - Test 20 (06 Apr 2024)}, paper=1, qno=9, type=single, subtopic={Area between curves}, answer={C}, answer_src=key, difficulty=4, diff_parts={K1 S2 I1 C1 F0}, diff_note={concepts: areas bounded by curves matched  values | steps: find intersection points; integrate piecewise; areas of three regions; match}]
Match the following:\par
{\footnotesize\setlength{\tabcolsep}{2pt}\noindent\begin{tabular}{|c|L{4.6cm}|c|L{1.0cm}|}
\hline
 & \multicolumn{1}{c|}{List-I} & & \multicolumn{1}{c|}{List-II}\\ \hline
(P) & Area bounded by $y=x^2+2$ and $y=2\,|x|-\cos\pi x$ is equal to $A$, then $3A$ equals & (1) & 4\\ \hline
(Q) & The value of \resizebox{\linewidth}{!}{$2\cot\left(\cot^{-1}3+\cot^{-1}7+\cot^{-1}13+\cot^{-1}21\right)$} & (2) & $\dfrac83$\\ \hline
(R) & Tangent and normal at the ends A \& C of focal chord AC of parabola $y^2=4x$ intersect at B and D, then minimum area of ABCD is & (3) & 3\\ \hline
(S) & Number of integral values of ordered pair $(\alpha,\beta)$ for which the area common to $x^2+y^2-2\alpha^2x-2\beta^2y+c=0$ and its image in $x+y=1$ is maximum is & (4) & 8\\ \hline
\end{tabular}}\par
Code:
\begin{qoptions}
\opt{A}{P $\rightarrow$ 4;Q $\rightarrow$ 3;R $\rightarrow$ 2;S $\rightarrow$ 3}
\opt{B}{P $\rightarrow$ 1;Q $\rightarrow$ 2;R $\rightarrow$ 1;S $\rightarrow$ 3}
\opt{C}{P $\rightarrow$ 4;Q $\rightarrow$ 3;R $\rightarrow$ 4;S $\rightarrow$ 1}
\opt{D}{P $\rightarrow$ 1;Q $\rightarrow$ 3;R $\rightarrow$ 4;S $\rightarrow$ 2}
\end{qoptions}
\end{bankq}
\begin{bankq}[id=SA21-P1-MATH-S1-Q05, source={Score Advanced - Test 21 (08 Apr 2024)}, paper=1, qno=5, type=multi, subtopic={Area between curves}, answer={B,C}, answer_src=key, difficulty=5, diff_parts={K1 S2 I1 C1 F1}, diff_note={concepts: area enclosed by two parabolas  axes  the shortest segment between them | steps: parabolas y^2 = x - 1 and x = y - 1 symmetric about the line y = x?; shortest distance along their common normal; area between curves and axes by integration; compare with expression}]
Consider the two parabolas $y^2=x-1$ and $x=y-1$. Let $A$ be the area between the coordinate axes, the two parabolas and the line which makes the shortest distance between two parabolas and $A^1$ be the area between $x=0$, $x^2=y-1$, $x=y$ and the line of shortest distance between $y^2=x-1$ and $x^2=y-1$ then which of the following is/are true?
\begin{qoptions}
\opt{A}{$A=A^1$}
\opt{B}{$A=2A^1$}
\opt{C}{The shortest distance between the parabolas is $\dfrac{3}{2\sqrt2}$}
\opt{D}{The shortest distance between the parabolas is $\dfrac12$}
\end{qoptions}
\end{bankq}
\begin{bankq}[id=SA22-P2-MATH-S1-Q05, source={Score Advanced - Test 22 (10 Apr 2024)}, paper=2, qno=5, type=single, subtopic={Area between curves}, answer={C}, answer_src=key, difficulty=5, diff_parts={K1 S2 I2 C1 F0}, diff_note={concepts: area under the inverse function of a cubic | steps: area of inverse between x = -3 and 5 using the integral of f over corresponding y limits and the rectangle; limits y from f^-1(-3) to f^-1(5); f(1)=5 f(-1) = -3; area = | I: area under the inverse function equals a rectangle minus the area of the original}]
Area enclosed by y = g(x), x = $-3$, x = 5 and x-axis where g(x) is the inverse function of $f(x)=x^3+3x+1$ is \_\_\_\_\_
\begin{qoptions}
\opt{A}{$\dfrac52$}
\opt{B}{$\dfrac72$}
\opt{C}{$\dfrac92$}
\opt{D}{$\dfrac32$}
\end{qoptions}
\end{bankq}
\begin{bankq}[id=SA24-P1-MATH-S3-Q01, source={Score Advanced - Test 24 (17 Apr 2024)}, paper=1, qno=1, type=numerical, subtopic={Area between curves}, answer={2}, answer_src=key, difficulty=6, diff_parts={K1 S3 I1 C1 F1}, diff_note={concepts: area of a region defined by inequalities  its weighted measure | steps: y >= 0, y <= x y? so for x >= 1; y <= 2 - x; sketch region as polygon pieces; area by integration; asked measure}]
In the coordinate plane, the region M consists of all the points (x, y) satisfying the inequalities $y\geq0$, $y\leq x\,y$ and $y\leq2-x$ simultaneously. The region N, which varies with the parameter t, consists of all the points (x, y) satisfying the inequalities $t\leq x\leq t+1$ and $0\leq t\leq1$ simultaneously. The area of M$\cap$N is $-k_1t^2+k_2t+\dfrac12$, then $k_1+k_2$ is \_\_\_\_\_\_ .
\end{bankq}
\begin{bankq}[id=SA24-P2-MATH-S1-Q01, source={Score Advanced - Test 24 (17 Apr 2024)}, paper=2, qno=1, type=single, subtopic={Area between curves}, answer={D}, answer_src=key, difficulty=4, diff_parts={K1 S2 I1 C1 F0}, diff_note={concepts: area bounded by a curve  its inverse region | steps: region bounded by sqrt(1+x^2) the line y=x y axis and y=-x+a; inverse function symmetric about y=x; area relations between the original and inverse regions; area of the inverse region in terms of k}]
If the area bounded by f (x) = $\sqrt{1+x^2}$ (x $>$ 0) the line y = x, y axis and y = $-$ x + a (a $>$1) is k, then area bounded by the graph of $f^{-1}$(x), the line y = x between the lines y = $-$ x + 1 and y = $-$ x + a is
\begin{qoptions}
\opt{A}{$\dfrac{k-1}{4}$}
\opt{B}{$\dfrac{2k-1}{2}$}
\opt{C}{k}
\opt{D}{$\dfrac{4k-1}{4}$}
\end{qoptions}
\end{bankq}
\begin{bankq}[id=SA25-P1-MATH-S2-Q05, source={Score Advanced - Test 25 (21 Apr 2024)}, paper=1, qno=5, type=multi, subtopic={Area between curves}, answer={A,C}, answer_src=key, difficulty=7, diff_parts={K1 S3 I2 C1 F1}, diff_note={concepts: areas of nested diamonds with recursive centres and sizes | steps: abs(x-a) + abs(y) = b is a square of diagonal 2b with area 2 b^2; sequences a_i b_i; areas form a geometric series 2 b_i^2; total area and overlap of consecutive diamonds; check statements | I: area of the diamond is 2 b^2}, also_in={Score Advanced - Test 26 (21 Apr 2024) P1 Q5}]
If $A_i$ is the area bounded by $|x-a_i|+|y|=b_i$, $i\in N$, where $a_{i+1}=a_i+\dfrac32b_i$ \& $b_{i+1}=\dfrac{b_i}{2}$, $a_1=0$, $b_1=32$, then
\begin{qoptions}
\opt{A}{$A_3=128$}
\opt{B}{$A_3=256$}
\opt{C}{$\displaystyle\lim_{n\to\infty}\sum_{i=1}^nA_i=\dfrac83(32)^2$}
\opt{D}{$\displaystyle\lim_{n\to\infty}\sum_{i=1}^nA_i=\dfrac43(16)^2$}
\end{qoptions}
\end{bankq}
\begin{bankq}[id=SA28-P1-MATH-S2-Q07, source={Score Advanced - Test 28 (27 Apr 2024)}, paper=1, qno=7, type=numerical, subtopic={Area between curves}, answer={17}, answer_src=key, difficulty=7, diff_parts={K1 S3 I2 C1 F1}, diff_note={concepts: area under the maximum  three quadratic functions on [0 1] | steps: compare x^2, (1-x)^2, 2x(1-x) on [0,1]; intersection points x = 1/3 and 2/3 (where 2x(1-x) overtakes); f(x) = (1-x)^2 on [0,1/3], 2x(1-x) on [1/3,2/3], x^2 on [2/3,1]; integrate three pieces; area lambda | I: identify the maximising function piecewise}]
Let f(x) = max\{x$^2$, (1 $-$ x)$^2$, 2x(1 $-$ x)\}, where 0 $\le$ x $\le$ 1. If the area bounded by y = f(x), x-axis, x = 0 \& x = 1 is $\lambda$, then 27($\lambda$) is :-
\end{bankq}
\begin{bankq}[id=SA28-P1-MATH-S2-Q10, source={Score Advanced - Test 28 (27 Apr 2024)}, paper=1, qno=10, type=numerical, subtopic={Area between curves}, answer={36}, answer_src=key, difficulty=6, diff_parts={K1 S2 I2 C1 F1}, diff_note={concepts: areas S_n under x^2 (1-x)^n  a limit of a sum | steps: S_n = B(3, n+1) = 2/((n+1)(n+2)(n+3)); sum of reciprocals or telescoping of S_n; partial fractions; sum to infinity | I: area is a beta integral with a telescoping closed form}]
Let S$_\mathrm n$ be the area of the figure enclosed by a curve y = x$^2$(1 $-$ x)$^\mathrm n$ (0 $\le$ x $\le$ 1) and x-axis, if $\displaystyle\lim_{n\to\infty}\sum_{k=1}^{n}\mathrm S_\mathrm k=\dfrac1{\sqrt{\mathrm A}}$, then `A' is equal to
\end{bankq}
\begin{bankq}[id=SA30-P2-MATH-S1-Q02, source={Score Advanced - Test 30 (04 May 2024)}, paper=2, qno=2, type=multi, subtopic={Area between curves}, answer={A,D}, answer_src=key, difficulty=6, diff_parts={K1 S3 I1 C1 F1}, diff_note={concepts: area enclosed between a logarithmic curve its inverse  the x axis | steps: y = ln(x + e) and x = ln(1/y) = -ln y i.e. y = e^{-x}; intersection point; area between curves and the x axis by integrals; evaluate in terms of e; properties of A}, also_in={Score Advanced - Test 31 (04 May 2024) P2 Q2}]
If A is the area enclosed between the curves y = $\ell$n(x + e), x = $\ell$n$\left(\dfrac1y\right)$ and x-axis then :
\begin{qoptions}
\opt{A}{A is prime number}
\opt{B}{A is composite number}
\opt{C}{A is irrational number}
\opt{D}{A$^2$ is even integer}
\end{qoptions}
\end{bankq}
\begin{bankpassage}[id=SA32-P2-MATH-S3-G3-4, source={Score Advanced - Test 32 (08 May 2024)}, questions={3-4}]
\textbf{Question Stem}\par
T is the region of the plane x + y + z = 1 with x, y, z $>$0. S is the set of points (a, b, c) in T such that just two of the following three inequalities hold : $a\leqslant\dfrac12$, $b\leqslant\dfrac13$, $c\leqslant\dfrac16$.
\end{bankpassage}
\begin{bankq}[id=SA32-P2-MATH-S3-Q03, source={Score Advanced - Test 32 (08 May 2024)}, paper=2, qno=3, type=numerical, subtopic={Area between curves}, answer={0.86}, answer_src=key, difficulty=4, diff_parts={K2 S1 I1 C1 F1 adj-1}, diff_note={concepts: area, the triangular region x + y + z = 1 with positive coordinates | steps: triangle with vertices (1,0,0) (0,1,0) (0,0,1) is equilateral of side sqrt2; area = sqrt3/2 | I: the plane section is an equilateral triangle}, group=SA32-P2-MATH-S3-G3-4]
Area of the region T is
\end{bankq}
\begin{bankq}[id=SA32-P2-MATH-S3-Q04, source={Score Advanced - Test 32 (08 May 2024)}, paper=2, qno=4, type=numerical, subtopic={Area between curves}, answer={0.34}, answer_src=key, difficulty=5, diff_parts={K2 S2 I1 C1 F1 adj-1}, diff_note={concepts: area, the region where exactly two of three inequalities hold | steps: T area sqrt3/2; regions in barycentric coordinates: areas of a>1/2 etc. in proportion to (1-1/2)^2 and so on; exactly two hold: complement computations with inclusion exclusion; fraction of T times area}, group=SA32-P2-MATH-S3-G3-4]
Area of the region S is
\end{bankq}
\begin{bankq}[id=SA38-P2-MATH-S1-Q02, source={Score Advanced - Test 38 (13 Apr 2025)}, paper=2, qno=2, type=single, subtopic={Area between curves}, answer={C}, answer_src=key, difficulty=4, diff_parts={K2 S2 I1 C1 F0 adj-1}, diff_note={concepts: area of the region bounded by parabolas, a line | steps: intersections of y^2 = 4x and y^2 = 12 - 2x at x = 2; line 3y + sqrt8 x = 5 sqrt8 cuts the region; set up integral of the lower curve between 0 and the intersection and the line; evaluate area}]
Let S = $\bigl\{(\mathrm x,\mathrm y)\in\mathbb R\times\mathbb R:\mathrm x\geqslant0,$ $\mathrm y\geqslant0,$ $\mathrm y^2\leqslant4\mathrm x,$ $\mathrm y^2\leqslant12-2\mathrm x$ and $3\mathrm y+\sqrt8\,\mathrm x\leqslant5\sqrt8\bigr\}$. If the area of the region S is $17\sqrt2(\alpha)$, then $\alpha$ is equal to
\begin{qoptions}
\opt{A}{$\dfrac{17}2$}
\opt{B}{$\dfrac{17}4$}
\opt{C}{$\dfrac13$}
\opt{D}{$\dfrac{17}3$}
\end{qoptions}
\end{bankq}
\begin{bankq}[id=SA39-P1-MATH-S4-Q09, source={Score Advanced - Test 39 (20 Apr 2025)}, paper=1, qno=9, type=numerical, subtopic={Area between curves}, answer={5}, answer_src=key, difficulty=5, diff_parts={K2 S2 I1 C1 F1 adj-1}, diff_note={concepts: area between exponential, logarithmic curves inside a square | steps: region in [0,2]x[0,2] with ln x <= y <= e^x capped by y = 2; integrate e^x up to where e^x = 2 and ln x piece; area = a - b ln 2; identify a b}]
If the area bounded by x = 0, y = 0, x = 2, y = 2, y $\leqslant$ e$^\mathrm x$ and y $\geqslant$ ln x is (a $-$ b ln2) sq. unit, where a $\in$ N, b $\in$ N then value of $\left(\dfrac{\mathrm a+\mathrm b}2\right)$ is
\end{bankq}
\begin{bankq}[id=SA40-P1-MATH-S1-Q04, source={Score Advanced - Test 40 (27 Apr 2025)}, paper=1, qno=4, type=multi, subtopic={Area between curves}, answer={A,C,D}, answer_src=key, difficulty=6, diff_parts={K1 S3 I2 C1 F1 adj-1}, diff_note={concepts: areas of nested diamonds with recursive centres and sizes | steps: abs(x-a) + abs(y) = b is a square of diagonal 2b with area 2 b^2; sequences a_i b_i; areas form a geometric series 2 b_i^2; total area and overlap of consecutive diamonds; check statements | I: area of the diamond is 2 b^2}]
If A$_\mathrm i$ is the area bounded by $|$x $-$ a$_\mathrm i|$ + $|$y$|$ = b$_\mathrm i$, i $\in$ N where a$_{\mathrm i+1}$ = a$_\mathrm i$ + $\dfrac32$b$_\mathrm i$ and b$_{\mathrm i+1}$ = $\dfrac{\mathrm b_\mathrm i}2$, a$_1$ = 0, b$_1$= 32 then
\begin{qoptions}
\opt{A}{A$_3$ = 128}
\opt{B}{A$_3$ = 256}
\opt{C}{$\displaystyle\lim_{\mathrm n\to\infty}\sum_{\mathrm i=1}^{\mathrm n}\mathrm A_\mathrm i=\dfrac83(32)^2$}
\opt{D}{$\dfrac{\mathrm A_6}{\mathrm A_{10}}$ = 256}
\end{qoptions}
\end{bankq}

\section*{SRG (Enthusiast) papers}
\begin{bankq}[id=SRG02-P1-MATH-S1-Q04, source={SRG (Enthusiast) - Test 02 (04 Oct 2020)}, paper=1, qno=4, type=multi, subtopic={Area between curves}, answer={B,D}, answer_src=key, working={Printed key.}, difficulty=6, diff_parts={K1 S3 I1 C1 F1}, diff_note={concepts: values of m for which the region between a parabola and a line has a given area | steps: intersection points at x = 0 and x = 1 - m; area = (1 - m)^3/6; set equal to 9/2; solve for m; check both sign options | I: area between a parabola and a line from the roots | C: solving the cubic condition}]
For which of the following values of $m$, is the area of the region bounded by the curve $y=x-x^2$ and the line $y=mx$ equals $9/2$?
\begin{qoptions}
\opt{A}{$-4$}
\opt{B}{$-2$}
\opt{C}{2}
\opt{D}{4}
\end{qoptions}
\end{bankq}
\begin{bankq}[id=SRG02-P1-MATH-S1-Q05, source={SRG (Enthusiast) - Test 02 (04 Oct 2020)}, paper=1, qno=5, type=multi, subtopic={Area between curves}, answer={A,C}, answer_src=key, working={Printed key.}, difficulty=7, diff_parts={K1 S2 I2 C2 F1}, diff_note={concepts: area between two exponential and logarithmic type curves | steps: compare the curves y = e^x x and y = x e^x ... find intersections; evaluate areas above and below the x axis; check each statement | I: identify the curves from their equations | C: integrals for each statement}]
Let $C_1: y=ex\ln x$ and $C_2: y=\dfrac{\ln x}{ex}$ be two curves, then
\begin{qoptions}
\opt{A}{Area bounded by these curves above the $x$-axis is zero}
\opt{B}{Area bounded by these curves below the $x$-axis is $\dfrac{e^2-4}{3e}$}
\opt{C}{Area bounded by these curves below the $x$-axis is $\dfrac{e^2-5}{4e}$}
\opt{D}{None of these}
\end{qoptions}
\end{bankq}
\begin{bankq}[id=SRG02-P1-MATH-S2-Q13, source={SRG (Enthusiast) - Test 02 (04 Oct 2020)}, paper=1, qno=13, type=numerical, subtopic={Area between curves}, answer={3}, answer_src=key, working={Printed key.}, difficulty=5, diff_parts={K1 S2 I1 C1 F1}, diff_note={concepts: area between a parabola and a line | steps: intersection points 0 and 1/3... of y^2 = 9x and y = 3x; area = integral of (3 sqrt x - 3x); m/n; m + n | I: intersection points for the integral | C: integral and reduction}]
The area of the figure bounded by $y^2=9x$ and $y=3x$ is $m/n$ units where $m$ and $n$ are coprime positive integers find $m+n$
\end{bankq}
\begin{bankq}[id=SRG02-P1-MATH-S2-Q14, source={SRG (Enthusiast) - Test 02 (04 Oct 2020)}, paper=1, qno=14, type=numerical, subtopic={Area between curves}, answer={4}, answer_src=key, working={Printed key.}, difficulty=5, diff_parts={K1 S2 I1 C1 F1}, diff_note={concepts: area between two modulus curves | steps: sketch y = abs(x) - 2 and y = 1 - abs(x - 1); find the intersection points; compute the enclosed region as a polygon | I: sketch both modulus graphs | C: polygon area}]
Area bounded by the curve $y=|x|-2$ and $y=1-|x-1|$ is equal to
\end{bankq}
\begin{bankq}[id=SRG02-P1-MATH-S2-Q15, source={SRG (Enthusiast) - Test 02 (04 Oct 2020)}, paper=1, qno=15, type=numerical, subtopic={Area between curves}, answer={9}, answer_src=key, working={Printed key.}, difficulty=5, diff_parts={K1 S2 I1 C1 F1}, diff_note={concepts: triangle formed by tangent normal and the x axis | steps: tangent and normal at (1,1) of y = x^2; intercepts on the x axis; area = half base times height; m/n | I: find the x intercepts of tangent and normal | C: area of the triangle}]
Area of the triangle formed by the $x$-axis and tangent and normal drawn to the curve $y=x^2$ at $(1,1)$ is $m/n$ where $m$ and $n$ are coprime positive integers then $m+n$ is equal to
\end{bankq}
\begin{bankq}[id=SRG02-P1-MATH-S2-Q16, source={SRG (Enthusiast) - Test 02 (04 Oct 2020)}, paper=1, qno=16, type=numerical, subtopic={Parametric regions}, answer={1}, answer_src=key, working={Printed key.}, difficulty=4, diff_parts={K1 S1 I1 C1 F1}, diff_note={concepts: area enclosed by a parametric curve | steps: x = (1 - t^2)/(1 + t^2), y = 2t/(1 + t^2) is the unit circle; area pi | I: recognise the parametrisation of the unit circle}]
Area enclosed by the curve $y=f(x)$ that is being defined parametrically as $x=\dfrac{1-t^2}{1+t^2}$, $y=\dfrac{2t}{1+t^2}$ is equal to $k\pi$ sq. units then $k$ is equal to
\end{bankq}
\begin{bankq}[id=SRG02-P1-MATH-S2-Q17, source={SRG (Enthusiast) - Test 02 (04 Oct 2020)}, paper=1, qno=17, type=numerical, subtopic={Area between curves}, answer={7}, answer_src=key, working={Printed key.}, difficulty=6, diff_parts={K2 S2 I1 C1 F1}, diff_note={concepts: area bounded by a normal, the axis and a parabola | steps: normal at (1,2) to y^2 = 4x: y = -x + 3; meets the x axis at x = 3; area between curve and normal from the vertex; m/n | I: region with the normal and the curve | C: integral of the region}]
The area bounded by the normal at $(1,2)$ on the curve, $y^2=4x$, $x$-axis and the curve is given by $m/n$ sq. units where $m$ and $n$ are coprime integers find $m-n$
\end{bankq}
\begin{bankq}[id=SRG02-P1-MATH-S2-Q18, source={SRG (Enthusiast) - Test 02 (04 Oct 2020)}, paper=1, qno=18, type=numerical, subtopic={Area between curves}, answer={3}, answer_src=key, working={Printed key.}, difficulty=7, diff_parts={K1 S2 I2 C2 F1}, diff_note={concepts: area bounded by three curves in the first quadrant | steps: identify the curves y = sqrt x, y = x ... and the x axis; intersection points; area as an integral; express as 2(m - n) form | I: the region is split at the intersection of the curves | C: integrals with radicals}]
The area of the figure bounded by $y=\sin x$, $y=\cos x$ and $x$-axis in the first quadrant is equal to $2\left(\sqrt m-n\right)$ where $m$ and $n$ are positive integers then find $m+n$
\end{bankq}
\begin{bankq}[id=SRG13-P1-MATH-S1-Q05, source={SRG (Enthusiast) - Test 13 (10 Oct 2021)}, paper=1, qno=5, type=single, subtopic={Area between curves}, answer={C}, answer_src=key, difficulty=5, diff_parts={K1 S2 I2 C2 F0 adj-1}, diff_note={concepts: area of intersection of a square and a circle | steps: find intersection points of the circle x squared + y squared = 4/3 with the unit square; split area into a triangle and a sector; add; match the form | I: splitting the region into a triangle and a sector | C: computing the sector area and angle}]
Consider the following regions in the plane

$R_1=\{(x,y):0\le x\le1\text{ and }0\le y\le1\}$ \quad $R_2=\{(x,y):x^2+y^2\le4/3\}$

The area of the region $R_1\cap R_2$ can be expressed as $\dfrac{a\sqrt3+b\pi}{9}$, where $a$ and $b$ are integers. Then the value of $(a+b)$ equals
\begin{qoptions}
\opt{A}{2}
\opt{B}{3}
\opt{C}{4}
\opt{D}{5}
\end{qoptions}
\end{bankq}
\begin{bankq}[id=SRG13-P1-MATH-S2-Q12, source={SRG (Enthusiast) - Test 13 (10 Oct 2021)}, paper=1, qno=12, type=multi, subtopic={Area between curves}, answer={A,D}, answer_src=key, difficulty=4, diff_parts={K2 S1 I1 C1 F1 adj-1}, diff_note={concepts: area under y = x^(1/3)-type curve, equivalent integrals | steps: area equals integral along x of the curve or along y of the inverse; evaluate each option | I: seeing the equivalent forms of the same area | C: checking four integrals}]
Area of the region bounded by the curve $y=\tan x$ and lines $y=0$ and $x=1$ is equal to
\begin{qoptions}
\opt{A}{$\displaystyle\int_0^1\tan(1-x)\,dx$}
\opt{B}{$\displaystyle\int_0^{\tan1}\tan^{-1}y\,dy$}
\opt{C}{$\displaystyle\int_0^1\tan^{-1}x\,dx$}
\opt{D}{$\tan1-\displaystyle\int_0^{\tan1}\tan^{-1}x\,dx$}
\end{qoptions}
\end{bankq}
\begin{bankq}[id=SRG13-P2-MATH-S3-Q11, source={SRG (Enthusiast) - Test 13 (10 Oct 2021)}, paper=2, qno=11, type=numerical, subtopic={Area between curves}, answer={5}, answer_src=key, difficulty=4, diff_parts={K1 S2 I1 C1 F1 adj-1}, diff_note={concepts: area between curves with exponential and logarithmic terms | steps: determine the region bounded by the curves; integrate over x; compare with the given form; compute a + b + c | I: identifying the boundaries of the region | C: evaluating the integral and matching the form}]
The area of the region bounded by the curves $2^{|x|}|y|+2^{|x|-1}\le1,\ |x|\le\dfrac12$ and $|y|\le\dfrac12$ is $\dfrac{a}{\log2}\left(1-\dfrac1{\sqrt b}\right)+c$ then $a+b+c$ is $(a,b,c\in N)$
\end{bankq}
\begin{bankq}[id=SRG13-P2-MATH-S3-Q12, source={SRG (Enthusiast) - Test 13 (10 Oct 2021)}, paper=2, qno=12, type=numerical, subtopic={Area between curves}, answer={2}, answer_src=key, difficulty=4, diff_parts={K1 S2 I1 C1 F1 adj-1}, diff_note={concepts: area between a parabola and two tangents from the origin | steps: tangent points x = plus minus 1; area = integral from -1 to 1 of difference; 3A | I: finding the tangent points from the origin | C: integrating the difference}]
The area bounded by $y=x^2+1$ and the tangents to it drawn from the origin is $A$ then $3A$ is
\end{bankq}
\begin{bankq}[id=SRG13-P2-MATH-S3-Q14, source={SRG (Enthusiast) - Test 13 (10 Oct 2021)}, paper=2, qno=14, type=numerical, subtopic={Area between curves}, answer={13}, answer_src=key, difficulty=6, diff_parts={K2 S2 I1 C2 F1 adj-1}, diff_note={concepts: area of region bounded by line, circle and parabola | steps: find intersection points; region splits; integrate; express as a b - 16 | I: finding the region from the three curves | C: integration of areas}]
The area given by $x+y\le6,\ x^2+y^2\le6y$ and $y^2\le8x$ is $\dfrac ab\pi-\dfrac16$ then $a+b$ is (given $a,b\in N$ $a$, $b$ are relatively prime)
\end{bankq}
\begin{bankq}[id=SRG14-P1-MATH-S1-Q01, source={SRG (Enthusiast) - Test 14 (21 Nov 2021)}, paper=1, qno=1, type=multi, subtopic={Area between curves}, answer={A,B,D}, answer_src=key, difficulty=5, diff_parts={K1 S3 I1 C1 F1 adj-1}, diff_note={concepts: area bounded by a parabola type curve and a line | steps: curves y = sqrt(-x) and x = -sqrt(y) i.e. same curve; line y = 1/4; area as integral of the difference; express as (1/3)(m/n) squared; test coprime conditions | I: recognising the two curves coincide | C: integration and matching the form}]
If area bounded by the curves $y\le\sqrt{-x},\ x\ge-\sqrt y$ and $y\ge\dfrac14$ is $\dfrac13\left(\dfrac mn\right)^2$, where $m$ \& $n$ are coprime numbers, then which of the following is (are) true-
\begin{qoptions}
\opt{A}{$|m-n|=1$}
\opt{B}{$m$ is prime \& $n$ is even}
\opt{C}{$n$ is prime number and $m$ is even number}
\opt{D}{$m+n$ is an odd number}
\end{qoptions}
\end{bankq}
\begin{bankq}[id=SRG14-P1-MATH-S2-Q09, source={SRG (Enthusiast) - Test 14 (21 Nov 2021)}, paper=1, qno=9, type=matching, subtopic={Area between curves}, answer={A$\to$T; B$\to$S; C$\to$P; D$\to$Q}, answer_src=key, difficulty=5, diff_parts={K3 S1 I1 C1 F1 adj-1}, diff_note={concepts: areas of curves and regions, minimum area, region defined by modulus | steps: area of the inverse sine curve; region of | x+y | + | x-y | is an area difference of squares; area between parabola and line minimised over parameter; match four areas | I: converting each region to a standard shape | C: evaluating four areas}]
Observe the following column:\par\nobreak
\begin{center}\begin{tabular}{|L{0.6cm}|L{8.4cm}|L{0.8cm}|L{2cm}|}\hline
\multicolumn{2}{|c|}{\textbf{Column I}} & \multicolumn{2}{c|}{\textbf{Column II}}\\\hline
(A) & Area enclosed by the curve $\left(y-\sin^{-1}x\right)^2=x-x^2$ is equal to & (P) & $0$\\\hline
(B) & The area of the region represented by the expression $\sqrt2\le|x+y|+|x-y|\le2\sqrt2$ is equal to & (Q) & $1$\\\hline
(C) & If the area bounded by $y=x^2-3$ and the line $y=ax+2$ attains its minimum value then the parameter `$a$' is equal to & (R) & $\pi$\\\hline
(D) & If $k$ is a positive number and the area of the region bounded by the curves $y=x-kx^2$ and $ky=x^2$ attains its maximum value then $k$ is equal to & (S) & $6$\\\hline
 & & (T) & $\pi/4$\\\hline
\end{tabular}\end{center}
\end{bankq}
\begin{bankq}[id=SRG14-P1-MATH-S3-Q17, source={SRG (Enthusiast) - Test 14 (21 Nov 2021)}, paper=1, qno=17, type=numerical, subtopic={Area between curves}, answer={4}, answer_src=key, difficulty=7, diff_parts={K2 S2 I2 C2 F1 adj-1}, diff_note={concepts: region defined by greatest integers, area | steps: 2[y] = 3[x] gives pairs of integer values [x] even... list each integer pair in the ranges; each pair gives a unit rectangle; count cells and add area | I: finding the integer pairs satisfying 2[y] = 3[x] within the ranges | C: counting rectangles of area 1 for each pair}]
$P(x,y)$ is a point, which moves in the $xy$ plane such that $2[y]=3[x]$, where $[\,.\,]$ denotes the greatest integer function and $-2\le x\le5$ and $-3\le y\le6$. The area of the region containing the point $P(x,y)$ is equal to
\end{bankq}
\begin{bankq}[id=SRG16-P1-MATH-S1-Q02, source={SRG (Enthusiast) - Test 16 (12 Jan 2022)}, paper=1, qno=2, type=single, subtopic={Area between curves}, answer={A}, answer_src=solved, working={The maximum area occurs for a = 2: tent and V cross at abs(x - 1) = 1/3 and the area is 1/3.}, difficulty=7, diff_parts={K2 S3 I2 C2 F0 adj-1}, diff_note={concepts: area between a tent function and a V function, maximisation over a parameter | steps: sketch both curves; intersection points depend on a; area as function of a; maximum at a = 2 by symmetry; area 1/3 | I: seeing the area is maximal when the vertex is under the apex | C: computing the area for general position}]
Let $a\in[0,4]$. The maximum area bounded by the curves $y=1-|x-1|$ \& $y=|2x-a|$ is
\begin{qoptions}
\opt{A}{1/3}
\opt{B}{1/4}
\opt{C}{1/2}
\opt{D}{None of these}
\end{qoptions}
\end{bankq}
\begin{bankq}[id=SRG16-P1-MATH-S1-Q05, source={SRG (Enthusiast) - Test 16 (12 Jan 2022)}, paper=1, qno=5, type=single, subtopic={Area between curves}, answer={C}, answer_src=solved, working={4 f cubed f prime = x squared f cubed gives f prime = x squared/4 so f = x cubed/12; area from -1 to 4 is (1/12)(1/4 + 64) = 257/48.}, difficulty=4, diff_parts={K2 S2 I1 C1 F0 adj-1}, diff_note={concepts: integral equation, area under cubic curve | steps: differentiate to f' = x squared/4; integrate; f(0) = 0; area from -1 to 4 using absolute value | I: dividing by f cubed after differentiating | C: area over two sub intervals}]
Let $f:[0,\infty)\to R$ be a continuous and strictly increasing function such that $f^4(x)=\displaystyle\int_0^x t^2f^3(t)\,dt\ \forall x>0$. The area enclosed by $y=f(x)$, the x-axis and the ordinates at $x=-1$ and $x=4$ is
\begin{qoptions}
\opt{A}{$\dfrac{17}{48}$}
\opt{B}{$\dfrac{65}{48}$}
\opt{C}{$\dfrac{257}{48}$}
\opt{D}{$\dfrac{235}{48}$}
\end{qoptions}
\end{bankq}
\begin{bankq}[id=SRG16-P2-MATH-S3-Q15, source={SRG (Enthusiast) - Test 16 (12 Jan 2022)}, paper=2, qno=15, type=numerical, subtopic={Area between curves}, answer={41}, answer_src=solved, working={Identity in t gives f(-1) = f(1) = 3/4, f(2) = 0, so f = 1 - x squared/4; V = (0,1), A = (2,0), B = (-8,-15); area = (1/4)(10) cubed/6 = 41.67, greatest integer 41.}, difficulty=7, diff_parts={K2 S2 I2 C2 F1 adj-1}, diff_note={concepts: quadratic from a functional identity, area under a chord of a parabola | steps: identity in t gives f(-1) = f(1) = 3/4 and f(2) = 0; f = 1 - x squared/4; vertex and chord end points; area (1/4)(10) cubed/6 = 41.67 | I: extracting f from the matrix equation | C: area under the chord and parabola}]
If $\begin{bmatrix}4a^2&4a&1\\4b^2&4b&1\\4c^2&4c&1\end{bmatrix}\begin{bmatrix}f(-1)\\f(1)\\f(2)\end{bmatrix}=\begin{bmatrix}3a^2+3a\\3b^2+3b\\3c^2+3c\end{bmatrix}$, $f(x)$ is a quadratic function and its maximum value occurs at a point V. A is a point of intersection of $y=f(x)$ with x-axis and point B is such that chord AB subtends a right angle at V. Let A be the area enclosed by $f(x)$ and chord AB. Find $[A]$ (where $[\cdot]$ denotes the greatest integer function)
\end{bankq}
\begin{bankq}[id=SRG38-P1-MATH-S1-Q06, source={SRG (Enthusiast) - Test 38 (15 Sep 2024)}, paper=1, qno=6, type=multi, subtopic={Area between curves}, answer={A,B,D}, answer_src=key, working={A: the integral reduces to the integral of x squared cos x over 0 to pi/2, pi squared/4-2. B: area 2 times the integral of x to the 3/2 is 4/5. C: the area between f and its inverse is 1/2, not 1/4. D: the normal at (1,2) cuts the parabola again at (9,-6), enclosing 64/3.}, difficulty=7, diff_parts={K2 S4 I1 C1 F1 adj-1}, diff_note={concepts: integral of odd and even parts, areas of curves | steps: reduce the integral by the property; integral of x squared cos x; area between branches; area between f and inverse; normal to the parabola; points of intersection; area of the segment; check each | I: using inverse symmetry | C: several areas}]
Which of the following is true
\begin{qoptions}
\opt{A}{$\displaystyle\int_{-\pi/2}^{\pi/2}\dfrac{x^2\cos x}{1+e^x}\,dx=\dfrac{\pi^2}{4}-2$}
\opt{B}{Area bounded by two branches of curve $(y-x)^2=x^3$ and $x=1$ is equal to $\dfrac45$}
\opt{C}{If $f(x)=\dfrac{x^3+x}{2}$, $f:R\to R$ \& $g(x)=f^{-1}(x)$ then area bounded between $f(x)$ and $g(x)$ is equal to $\dfrac14$}
\opt{D}{Area bounded by $y^2=4x$ and normal to $y^2=4x$ with slope $-1$ is equal to $\dfrac{64}{3}$}
\end{qoptions}
\end{bankq}
\begin{bankq}[id=SRG38-P1-MATH-S1-Q09, source={SRG (Enthusiast) - Test 38 (15 Sep 2024)}, paper=1, qno=9, type=multi, subtopic={Area between curves}, answer={A,B,C}, answer_src=key, working={On [0,10] f=max(x,10-x), area 75. g is the distance to the nearest integer, area 2.5. The area between 1 and g is 7.5 and between f and g is 75-2.5=72.5.}, difficulty=6, diff_parts={K2 S3 I1 C1 F1 adj-1}, diff_note={concepts: graphs of maximum and minimum of modulus functions, area | steps: form of f; area of f; form of g; area from y equal 1; area between f and g | I: shapes of f and g | C: triangle areas}]
If $f(x)=\max\{|x|,|x-1|,|x-2|,\ldots,|x-10|\}$ and $g(x)=\min\{|x|,|x-1|,|x-2|,\ldots|x-10|\}$ then
\begin{qoptions}
\opt{A}{Area bounded by $y=f(x)$, $x$-axis, $x=0$ and $x=10$ is 75 unit}
\opt{B}{Area bounded by $y=g(x)$, $y=1$, $x=0$ and $x=10$ is 7.5}
\opt{C}{Area bounded by $y=f(x)$, $y=g(x)$, $x=0$ and $x=10$ is 72.5}
\opt{D}{Area bounded by $y=f(x)$, $y=g(x)$, $x=0$ and $x=10$ is 70}
\end{qoptions}
\end{bankq}
\begin{bankq}[id=SRG38-P1-MATH-S2-Q06, source={SRG (Enthusiast) - Test 38 (15 Sep 2024)}, paper=1, qno=6, type=numerical, subtopic={Area between curves}, answer={0}, answer_src=key, working={f=ln squared x-1/6 and the tangent at e has slope 2/e. The area is the integral of the difference on [1,e], e+1/e-3=0.086, so [A]=0.}, difficulty=6, diff_parts={K2 S3 I1 C1 F1 adj-1}, diff_note={concepts: integral equation for a function, area of a region between a curve and its tangent | steps: constant from the equation; the function; tangent at e; area integral; evaluate; floor | I: constant integral trick | C: integrals of ln squared x}]
Let $f(x)$, $(x\ge1)$ be a differentiable function satisfying $f(x)=\ell n^2x-\displaystyle\int_1^{e}\dfrac{f(t)}{t}\,dt$, then area bounded by tangent line of $y=f(x)$ at $x=e$, the curve $y=f(x)$ and the line $x=1$ is $A$, then $[A]$ is equal to (where $[\,.\,]$ denotes greatest integer function)
\end{bankq}
\begin{bankq}[id=SRG38-P1-MATH-S2-Q10, source={SRG (Enthusiast) - Test 38 (15 Sep 2024)}, paper=1, qno=10, type=numerical, subtopic={Area between curves}, answer={7}, answer_src=key, working={g=x-c with c the integral of f, so f=1+c x squared and c=1+c/3, c=3/2. The area between f=1+1.5x squared and g=x-1.5 on 0 to 2 is 4-2+5=7.}, difficulty=6, diff_parts={K2 S3 I1 C1 F1 adj-1}, diff_note={concepts: integral equations, area between curves | steps: g in terms of a constant; f in terms of the constant; solve; difference of the functions; integrate | I: treating the integral as a constant | C: integration}]
$f(x)$ and $g(x)$ are two differentiable function defined on $R$, satisfying $f(x)=\dfrac{x^3}{2}+1-x\displaystyle\int_0^{x}g(t)\,dt$ and $g(x)=x-\displaystyle\int_0^{1}f(t)\,dt$, then the area bounded by $y=f(x)$, $y=g(x)$, $x=0$ and $x=2$ is equal to
\end{bankq}

\section*{Pre-Score (2023-24) papers}
\begin{bankq}[id=PS07-P1-MATH-S1-Q05, source={Pre-Score (2023-24) - Test 07 (29 Nov 2023)}, paper=1, qno=5, type=multi, subtopic={Area between curves}, answer={A}, answer_src=key, working={Printed key.}, difficulty=9, diff_parts={K1 S4 I2 C2 F1}, diff_note={concepts: equal area regions in the first quadrant bounded by an ellipse | steps: area of region under 2y = x inside the ellipse; area of region with y = mx; equate; use ellipse area in sectors with inverse sine; solve for m; area of region R from 0 to the intersection; locate the intersection of 2y = x with the ellipse; area of R prime by sector method; equate the areas; solve m; compare with options | I: sector area of an ellipse by scaling a circle | C: equating the areas}]
Let $R$ be the region in the first quadrant bounded by the $x$-axis, the line $2y=x$, and the ellipse $\dfrac{x^2}{9}+y^2=1$. Let $R'$ be the region in the first quadrant bounded by the $y$-axis, the line $y=mx$, and the given ellipse, then the value of $m$ such that $R$ and $R'$ have the same area is
\begin{qoptions}
\opt{A}{$\dfrac29$}
\opt{B}{$\dfrac3{11}$}
\opt{C}{$\dfrac25$}
\opt{D}{$\dfrac27$}
\end{qoptions}
\end{bankq}
\begin{bankq}[id=PS07-P1-MATH-S2-Q10, source={Pre-Score (2023-24) - Test 07 (29 Nov 2023)}, paper=1, qno=10, type=single, subtopic={Area between curves}, answer={C}, answer_src=key, working={Printed key.}, difficulty=5, diff_parts={K1 S2 I2 C1 F0}, diff_note={concepts: areas bounded by modulus inequalities with exponential functions | steps: solve for the boundary curve absolute y = (1 - 2 to the power (abs x - 1))/2 to the power abs x; area by symmetry; integrals of 2 to the power minus x; match four areas | I: use symmetry in x and y | C: four integrals}]
Match List-I with List-II :\par\nobreak
\begin{center}\begin{tabular}{|L{0.8cm}|L{6.3cm}|L{0.8cm}|L{4.2cm}|}\hline
\multicolumn{2}{|c|}{\textbf{List-I}} & \multicolumn{2}{c|}{\textbf{List-II}}\\\hline
(I) & The area bounded by $2^{|x|}|y|+2^{|x|-1}\le1$ is & (P) & $4-\dfrac{2}{\ell n2}$\\\hline
(II) & The area bounded by $2^{|x|}|y|+2^{|x|-1}\le1$, $|x|\le\dfrac12$ \& $|y|\le\dfrac12$ is equal to & (Q) & $\dfrac{2}{\ell n2}-2$\\\hline
(III) & The area bounded by $2^{|x|}|y|+2^{2|x|-1}\ge1$ \& $|x|+|y|\le1$ is equal to & (R) & $\dfrac{4\left(2^{1/3}-1\right)}{2^{1/3}\ell n2}-\dfrac23$\\\hline
(IV) & The area bounded by $2^{|x|}|y|+2^{|x|-1}\le1$ \& $|x|\le\dfrac13$ & (S) & $\dfrac{\left(\sqrt2-1\right)2\sqrt2}{\ell n2}-1$\\\hline
 & & (T) & $\dfrac{2}{\ell n2}+2$\\\hline
\end{tabular}\end{center}
\begin{qoptions}
\opt{A}{I $\to$ P; II $\to$ Q; III $\to$ R; IV $\to$ S}
\opt{B}{I $\to$ R; II $\to$ S; III $\to$ Q; IV $\to$ P}
\opt{C}{I $\to$ Q; II $\to$ S; III $\to$ P; IV $\to$ R}
\opt{D}{I $\to$ S; II $\to$ P; III $\to$ R; IV $\to$ Q}
\end{qoptions}
\end{bankq}
\begin{bankq}[id=PS07-P1-MATH-S3-Q04, source={Pre-Score (2023-24) - Test 07 (29 Nov 2023)}, paper=1, qno=4, type=numerical, subtopic={Area between curves}, answer={8}, answer_src=solved, working={Area under the curve from 0 to 3/2 is 39/8; the region under the line is a triangle of area 9m/8 and must be 39/16 so m = 13/6; 2p - 3q = 26 - 18 = 8.}, difficulty=6, diff_parts={K1 S3 I1 C1 F1}, diff_note={concepts: line through the origin bisecting an area under a parabola | steps: area under the curve from 0 to 1.5 is 39/8; triangle under the line has area 9m/8; set equal to half; solve m; reduce the fraction; compute 2p minus 3q | I: the region under the line is a triangle | C: area integral and fraction}]
If the line $y=mx$ bisects the area enclosed by the lines $x=0$, $y=0$, $x=\dfrac32$ and the curve $y=1+4x-x^2$ for $m=\dfrac pq$ ($p,q\in N$, $p$ \& $q$ are co-prime), then the value of $2p-3q=$
\end{bankq}
\begin{bankq}[id=PS07-P1-MATH-S3-Q06, source={Pre-Score (2023-24) - Test 07 (29 Nov 2023)}, paper=1, qno=6, type=numerical, subtopic={Area between curves}, answer={0.46 to 0.48}, answer_src=key, working={Printed key.}, difficulty=5, diff_parts={K1 S2 I1 C1 F1}, diff_note={concepts: area of the region defined by a sum of absolute values | steps: change variables u = 2x + 3y, v = 3x - 4y; region abs u + abs v at most 2 is a diamond of area 8; Jacobian is 17; area = 8/17 | I: linear change of variables converts the region to a diamond | C: Jacobian and area}]
Area of points satisfying $|2x+3y|+|3x-4y|\le2$ equals
\end{bankq}
\begin{bankq}[id=PS07-P2-MATH-S1-Q01, source={Pre-Score (2023-24) - Test 07 (29 Nov 2023)}, paper=2, qno=1, type=single, subtopic={Area between curves}, answer={C}, answer_src=key, working={Printed key.}, difficulty=4, diff_parts={K1 S2 I1 C1 F0}, diff_note={concepts: areas of regions defined by inequalities on circles and lines | steps: draw the three regions in the first quadrant; compute areas in terms of pi; compare | I: understand the regions from the inequalities | C: area of overlapping regions}]
Let $A_1$, $A_2$ and $A_3$ be the regions on $\mathbb R^2$ defined by\\
$A_1=\{(x,y):x\ge0,\ y\ge0,\ 2x+2y-x^2-y^2>1>x+y\}$,\\
$A_2=\{(x,y):x\ge0,\ y\ge0,\ x+y>1>x^2+y^2\}$,\\
$A_3=\{(x,y):x\ge0,\ y\ge0,\ x+y>1>x^3+y^3\}$\\
Denote by $|A_1|,|A_2|$ and $|A_3|$ the areas of the regions $A_1$, $A_2$ and $A_3$ respectively. Then
\begin{qoptions}
\opt{A}{$|A_1|>|A_2|>|A_3|$}
\opt{B}{$|A_1|>|A_3|>|A_2|$}
\opt{C}{$|A_1|=|A_2|<|A_3|$}
\opt{D}{$|A_1|=|A_3|>|A_2|$}
\end{qoptions}
\end{bankq}
\begin{bankq}[id=PS07-P2-MATH-S2-Q07, source={Pre-Score (2023-24) - Test 07 (29 Nov 2023)}, paper=2, qno=7, type=multi, subtopic={Area between curves}, answer={B,D}, answer_src=key, working={Printed key.}, difficulty=5, diff_parts={K1 S2 I1 C1 F1}, diff_note={concepts: area between the curves abs sin x and sin x over n half periods | steps: areas in each half period are twice the integral of sin over alternating intervals; sum for odd or even n; compare with formulas | I: sign of sin x alternates by half period | C: summing for odd and even n}]
The area bounded by the function $f(x)=|\cos x|$ and $g(x)=\sin x$ between the ordinates $x=0$ and $x=n\pi$, $n\in N$, is equal to
\begin{qoptions}
\opt{A}{$4n\sqrt2$ if $n$ is odd}
\opt{B}{$2n\sqrt2$ if $n$ is even}
\opt{C}{$2\sqrt2\,(n+1)$ if $n$ is odd}
\opt{D}{$2\sqrt2\,(n+1)-4$ if $n$ is odd}
\end{qoptions}
\end{bankq}
\begin{bankq}[id=PS07-P2-MATH-S2-Q08, source={Pre-Score (2023-24) - Test 07 (29 Nov 2023)}, paper=2, qno=8, type=multi, subtopic={Area between curves}, answer={A,B,C}, answer_src=key, working={Printed key.}, difficulty=5, diff_parts={K1 S2 I1 C1 F1}, diff_note={concepts: ordering of areas under four power-exponent curves | steps: compare the functions on each interval; larger base and exponent give larger values; compare areas A1 A2 A3 A4 | I: compare pointwise using the sign of the logarithm | C: reasoning about four curves}]
If the Area bounded by $y=(\sin x)^{\mathrm{cosec}\,x}$, $y=(\mathrm{cosec}\,x)^{\sin x}$, $y=(\sin x)^{\sin x}$ and $y=(\mathrm{cosec}\,x)^{\mathrm{cosec}\,x}$ between the ordinates $x=0$, $x=\dfrac\pi2$ denoted by $A_1$, $A_2$, $A_3$ and $A_4$ then
\begin{qoptions}
\opt{A}{$A_4$ is the greatest}
\opt{B}{$A_1$ is the least}
\opt{C}{$A_2>A_3$}
\opt{D}{$A_2<A_3$}
\end{qoptions}
\end{bankq}
\begin{bankq}[id=PS07-P2-MATH-S3-Q04, source={Pre-Score (2023-24) - Test 07 (29 Nov 2023)}, paper=2, qno=4, type=numerical, subtopic={Area between curves}, answer={3}, answer_src=key, working={Printed key.}, difficulty=5, diff_parts={K1 S2 I1 C1 F1}, diff_note={concepts: area between two exponential curves and a vertical line | steps: find the intersection and compare f and g; integrate g - f on the interval to x = 1; express as (2 - sqrt3) type plus e term | I: f and g involve absolute exponent so split at x = 1 | C: integrating exponentials and simplifying}]
Let the functions $f:R\to R$ and $g:R\to R$ be defined by $f(x)=e^{x-1}-e^{-|x-1|}$ and $g(x)=\dfrac12\left(e^{x-1}+e^{1-x}\right)$. Then the area of the region bounded by the curves $y=f(x)$, $y=g(x)$ and $x=1$ is $2-\sqrt\lambda$, then the value of $\lambda$ is
\end{bankq}
\begin{bankq}[id=PS07-P2-MATH-S3-Q05, source={Pre-Score (2023-24) - Test 07 (29 Nov 2023)}, paper=2, qno=5, type=numerical, subtopic={Area between curves}, answer={9}, answer_src=key, working={Printed key.}, difficulty=5, diff_parts={K1 S2 I1 C1 F1}, diff_note={concepts: ellipse through the endpoint of the latus rectum of a parabola | steps: endpoint of the latus rectum of y^2 = 4 lambda x is (lambda, 2 lambda); substitute into the ellipse; use the given relations among a b lambda; solve the requested quantity | I: choose the latus rectum point and apply the ellipse condition | C: algebra with the parameters}]
Let $a$, $b$ and $\lambda$ be positive real numbers. Suppose $P$ is an end point of the latus rectum of the parabola $y^2=4\lambda x$, and suppose the ellipse $\dfrac{x^2}{a^2}+\dfrac{y^2}{b^2}=1$ passes through the point $P$. If the tangents to the parabola and the ellipse at the point $P$ are perpendicular to each other, then the eccentricity of the ellipse is $e$. If $\left(e,\dfrac ab\right)$ is a point on the ellipse $15x^2+3y^2=K$, then $K$ is equal to \underline{\hspace{1.5cm}}
\end{bankq}
\begin{bankq}[id=PS07-P2-MATH-S3-Q07, source={Pre-Score (2023-24) - Test 07 (29 Nov 2023)}, paper=2, qno=7, type=numerical, subtopic={Area between curves}, answer={5}, answer_src=key, working={Printed key.}, difficulty=5, diff_parts={K1 S2 I1 C1 F1}, diff_note={concepts: area between a curve defined by an integral and the axes | steps: differentiate the integral definition to get f prime; find where f vanishes; integrate f to find the area as p/q | I: the function defined by two integrals simplifies on differentiation | C: integration and fraction}]
The area bounded by the curve $y=f(x)=\displaystyle\int_1^{1/x^3}\left(\dfrac1{1+t^4}+t^{-4/3}\right)dt-3\int_x^1\dfrac{t^8}{1+t^{12}}\,dt$ and the co-ordinate axes is equal to $\dfrac pq$ (where $p,q\in N$, G.C.D $=1$) then the value of $p+q$ is
\end{bankq}

\section*{JEE Advanced (2021 Series) papers}
\begin{bankq}[id=JA03-P1-MATH-S3-Q06, source={JEE Advanced (2021 Series) - Test 03 (01 Mar 2021)}, paper=1, qno=6, type=numerical, subtopic={Area between curves}, answer={9.00}, answer_src=key, difficulty=4, diff_parts={K1 S2 I1 C0 F1}, diff_note={concepts: area of a region defined with greatest integer functions | steps: for x in [n, n+1) the relation gives y in [2^n, 2^n + 1); each strip is a unit square for n = 0 to 8; total area 9}]
Area enclosed within the curve $[y]=2^{[x]}$, $x\in[0,9)$ (where $[\cdot]$ denotes the greatest integer function) is equal to
\end{bankq}
\begin{bankq}[id=JA03-P2-MATH-S3-Q04, source={JEE Advanced (2021 Series) - Test 03 (01 Mar 2021)}, paper=2, qno=4, type=numerical, subtopic={Area between curves}, answer={1.00}, answer_src=key, difficulty=6, diff_parts={K2 S2 I2 C0 F1}, diff_note={concepts: differentiating a functional equation to get a differential equation, area by improper integral | steps: set y = 0 in the equation and differentiate in y at 0 to get f'(x) = -f(x)^2 using f'(0) = -1; solve f = 1/(1 + x); curve y = 1/(1 + x)^2 and axes; integrate from 0 to infinity giving 1 | I: derivative with respect to y at y = 0 converts the functional equation to an ODE}]
Let $f:(-1,\infty)\to(0,\infty)$ be a differentiable function satisfying $f(x)\cdot f(yf(x))=f(x+y)\ \forall\ x,y$ (for which $f$ is well defined in given domain) and $f'(0)=-1$ and $f(0)=1$ then area of the region bounded by the curve $y=\dfrac{f(x)}{x+1}$ and the coordinate axes is
\end{bankq}
\begin{bankq}[id=JA05-P1-MATH-S1-Q03, source={JEE Advanced (2021 Series) - Test 05 (11 Apr 2021)}, paper=1, qno=3, type=single, subtopic={Area between curves}, answer={C}, answer_src=key, difficulty=3, diff_parts={K1 S2 I1 C0 F0}, diff_note={concepts: area under a periodic fractional part function | steps: 3 - x has the same fractional part as 1 - {x}; so f = min({x}, 1 - {x}) which is a triangle of area 1/4 per unit; 8 units from -4 to 4 give 2}]
Area bounded by $f(x)=\min(\{x-2\},\{3-x\})$, $x=4$, $x=-4$ and $x$-axis is equal to (where $\{.\}$ denotes fractional part)
\begin{qoptions}
\opt{A}{$1$}
\opt{B}{$3/2$}
\opt{C}{$2$}
\opt{D}{$4$}
\end{qoptions}
\end{bankq}
\begin{bankq}[id=JA05-P2-MATH-S3-Q02, source={JEE Advanced (2021 Series) - Test 05 (11 Apr 2021)}, paper=2, qno=2, type=numerical, subtopic={Area between curves}, answer={1}, answer_src=key, difficulty=5, diff_parts={K1 S2 I1 C1 F1}, diff_note={concepts: area between two parabolas as a function of a parameter with maximisation | steps: intersection of y = x - k x^2 and y = x^2/k at x = 0 and x = k/(k^2 + 1); area = (1/6)(k/(k^2 + 1))^3 (1 + 1/k)... simplify; maximise over k > 0; k = 1}]
The value of $k(k\in\mathbb R^+)$ so that area of the bounded region enclosed between the parabolas $y=x-kx^2$ and $y=\dfrac{x^2}{k}$ is maximum, is
\end{bankq}
\begin{bankq}[id=JA11-P2-MATH-S2-Q13, source={JEE Advanced (2021 Series) - Test 11 (20 May 2021)}, paper=2, qno=13, type=multi, subtopic={Area between curves}, answer={C,D}, answer_src=key, difficulty=4, diff_parts={K2 S2 I0 C0 F1}, diff_note={concepts: area between curves with symmetry, regions defined by modulus | steps: by symmetry take the first quadrant; A = 4 integral of ((1 - y^2) - (1 - y)) dy from 0 to 1 = 2/3; B = 2 integral of (x - x^2) dx = 1/3; ratio A/B = 2 which is less than 3 and 4}, also_in={JEE Advanced (2021 Series) - Test 15 (12 Jun 2021) P2 Q6; Nurture (General) - Test 01 (12 Jun 2021) P2 Q6}]
Area of region given by curves $1-y^2\ge|x|$ and $|x|+|y|\ge1$ is $A$ and area of region $\{(x,y):x^2\le y\le|x|\}$ is $B$. Then $\dfrac AB$ is less than
\begin{qoptions}
\opt{A}{$1$}
\opt{B}{$2$}
\opt{C}{$3$}
\opt{D}{$4$}
\end{qoptions}
\end{bankq}
\begin{bankq}[id=JA13-P2-MATH-S5-Q02, source={JEE Advanced (2021 Series) - Test 13 (31 May 2021)}, paper=2, qno=2, type=numerical, subtopic={Area between curves}, answer={1}, answer_src=key, working={Printed key.}, difficulty=5, diff_parts={K1 S2 I2 C0 F1}, diff_note={concepts: area enclosed by a quartic curve through polar coordinates | steps: put x = r cos t and y = r sin t to get r^2 = sin 2t/(1 - (1/2) sin^2 2t); area of one loop is (1/2) integral of r^2 dt = pi/4 using u = cos 2t; the curve has two loops so A = pi/2 = 1.57 and the integer part is 1 | I: polar substitution reduces the quartic to a rational integral in cos 2t}]
Let the area enclosed by the curve $x^4+y^4=2xy$ be $A$ square units, then the value of $[A]$ is equal to (where $[.]$ denotes the greatest integer function)
\end{bankq}
\begin{bankpassage}[id=JA14-P2-MATH-S2-G3-4, source={JEE Advanced (2021 Series) - Test 14 (06 Jun 2021)}, questions={3-4}]
Consider the area $S_n$ bounded by $x$-axis, $y=e^{-x}\sin x$ between $x=n\pi$ and $x=(n+1)\pi$, where $x\ge0$\\
On the basis of above information, answer the following questions :
\end{bankpassage}
\begin{bankq}[id=JA14-P2-MATH-S2-Q03, source={JEE Advanced (2021 Series) - Test 14 (06 Jun 2021)}, paper=2, qno=3, type=single, subtopic={Area between curves}, answer={B}, answer_src=key, working={Printed key.}, difficulty=4, diff_parts={K1 S2 I1 C1 F0}, diff_note={concepts: area bounded by e to the minus x sin x on the first arch | steps: S0 is integral of e^(-x) sin x from 0 to pi; use integration by parts formula; evaluate with limits; get half of one plus e to the minus pi | I: recognise the standard integral of e^(-x) sin x | C: evaluating at the limits and simplifying}, group=JA14-P2-MATH-S2-G3-4]
The value of $S_0$, is-
\begin{qoptions}
\opt{A}{$\dfrac12\left(1+e^{\pi}\right)$ sq. units}
\opt{B}{$\dfrac12\left(1+e^{-\pi}\right)$ sq. units}
\opt{C}{$\dfrac12\left(1-e^{-\pi}\right)$ sq. units}
\opt{D}{$\dfrac12\left(e^{\pi}-1\right)$ sq. units}
\end{qoptions}
\end{bankq}
\begin{bankq}[id=JA14-P2-MATH-S2-Q04, source={JEE Advanced (2021 Series) - Test 14 (06 Jun 2021)}, paper=2, qno=4, type=single, subtopic={Area between curves}, answer={B}, answer_src=key, working={Printed key.}, difficulty=4, diff_parts={K1 S2 I1 C1 F0}, diff_note={concepts: sum of successive arch areas as a geometric series | steps: areas S_n alternate and are in geometric ratio e to the minus pi; use S_0 as first term; common ratio e^(-pi); sum infinite series | I: successive areas form a geometric progression with ratio e^(-pi) | C: simplification of infinite series sum}, group=JA14-P2-MATH-S2-G3-4]
$\displaystyle\sum_{n=0}^{\infty}S_n$ is equal to-
\begin{qoptions}
\opt{A}{$\dfrac{1+e^{\pi}}{1-e^{-\pi}}$}
\opt{B}{$\dfrac12\dfrac{\left(1+e^{-\pi}\right)}{\left(1-e^{-\pi}\right)}$}
\opt{C}{$\dfrac1{2\left(1-e^{-\pi}\right)}$}
\opt{D}{none of these}
\end{qoptions}
\end{bankq}

\section*{TOAS papers}
\begin{bankq}[id=TOAS02-P1-MATH-S3-Q06, source={TOAS - Test 02 (13 Feb 2022)}, paper=1, qno=6, type=numerical, subtopic={Area between curves}, answer={9}, answer_src=solved, working={For $x\in[n,n+1)$: $[y]=2^n$ gives a unit square, $n=0,\dots,8$: area 9}, difficulty=3, diff_parts={K1 S1 I1 C0 F1}, diff_note={concepts: region given by greatest integers | steps: for each n the region is a unit square; nine of them | I: unit squares}]
Area enclosed within the curve $[y]=2^{[x]}$, $x\in[0,9)$ (where [$\bullet$] denotes the greatest integer function) is equal to
\end{bankq}
\begin{bankq}[id=TOAS06-P1-MATH-S2-Q08, source={TOAS - Test 06 (17 Apr 2022)}, paper=1, qno=8, type=multi, subtopic={Area between curves}, answer={B}, answer_src=key, difficulty=3, diff_parts={K1 S1 I1 C0 F1}, diff_note={concepts: area of a region with modulus | steps: region  | x}]
Let $S=\{(x,y)\mid x^2+y^2\geqslant4\}$ and $T=\{(x,y)\mid x^2-y^2\leqslant0\}$. Area of the region given by the set $T-S$ (in sq. units) is
\begin{qoptions}
\opt{A}{$\pi$}
\opt{B}{$2\pi$}
\opt{C}{$3\pi$}
\opt{D}{$4\pi$}
\end{qoptions}
\end{bankq}
\begin{bankq}[id=TOAS13-P1-MATH-S1-Q03, source={TOAS - Test 13 (24 Jul 2022)}, paper=1, qno=3, type=multi, subtopic={Area between curves}, answer={A,B,C,D}, answer_src=key, difficulty=3, diff_parts={K1 S2 I0 C0 F1}, diff_note={concepts: area between symmetric curves | steps: both curves are symmetric about x = 1; intersections at x = 0, 1, 2; area 2 times the integral on one side; check each option form}]
Area bounded by curves $y=\sqrt{|x-1|}$ and $\sqrt y=|x-1|$ is
\begin{qoptions}
\opt{A}{$2\displaystyle\int_1^2\left(\sqrt{x-1}-(x-1)^2\right)dx$}
\opt{B}{$2\displaystyle\int_0^1\left(\sqrt x-x^2\right)dx$}
\opt{C}{$2\displaystyle\int_1^2\left(\sqrt{2-x}-(2-x)^2\right)dx$}
\opt{D}{$\dfrac23$}
\end{qoptions}
\end{bankq}
\begin{bankq}[id=TOAS13-P1-MATH-S3-Q05, source={TOAS - Test 13 (24 Jul 2022)}, paper=1, qno=5, type=numerical, subtopic={Area between curves}, answer={1}, answer_src=solved, working={$\int_0^{1/2}f^{-1}=\tfrac12\cdot1-\int_0^1\dfrac{t}{1+t^2}dt=\tfrac12-\tfrac12\ln2$; $f^{-1}$ is odd so area $=2\left(\tfrac12-\tfrac12\ln2\right)=1-\ln2$, $A=1$}, difficulty=4, diff_parts={K1 S2 I1 C0 F1}, diff_note={concepts: area through the inverse function | steps: area is twice the integral of f\^{}-1 on 0 to 1/2; swap to the integral of f; evaluate}]
$f:[-1,1]\to\left[-\dfrac12,\dfrac12\right]$, $f(x)=\dfrac{x}{1+x^2}$. If the area bounded by $y=f^{-1}(x)$, X-axis between, $x=\dfrac{-1}{2}$ to $x=\dfrac12$ is $A-\ln2$, then the value of $A$ is equal to \underline{\hspace{1.5cm}}.
\end{bankq}
\begin{bankq}[id=TOAS13-P2-MATH-S1-Q14, source={TOAS - Test 13 (24 Jul 2022)}, paper=2, qno=14, type=single, subtopic={Area between curves}, answer={B}, answer_src=key, difficulty=1, diff_parts={K1 S1 I0 C0 F0}, diff_note={concepts: area between two parabolas and a line | steps: intersection at x = 2/3; integrate the difference to x = 1}]
The area bounded by the curves $f(x)=9x^2-9x+2$, $g(x)=9x^2-18x+8$ and $x=1$ is :
\begin{qoptions}
\opt{A}{$\dfrac13$}
\opt{B}{$\dfrac12$}
\opt{C}{$\dfrac23$}
\opt{D}{$\dfrac34$}
\end{qoptions}
\end{bankq}
\begin{bankq}[id=TOAS13-P2-MATH-S2-Q04, source={TOAS - Test 13 (24 Jul 2022)}, paper=2, qno=4, type=numerical, subtopic={Area between curves}, answer={2}, answer_src=solved, working={Reflection in $x+y=0$ maps $f$ to its inverse $g=f^{-1}$ ($f$ is odd). $\int_0^{2\pi}f^{-1}=2\pi\cdot2\pi-\int_0^{2\pi}(t+\sin t)dt=4\pi^2-2\pi^2=2\pi^2$, so $A/\pi^2=2$}, difficulty=4, diff_parts={K2 S1 I1 C0 F1}, diff_note={concepts: reflection about x + y = 0 gives the inverse function, area through the inverse | steps: g = f\^{}-1 because f is odd; area = 2 pi times 2 pi minus the integral of f}]
Let $y=g(x)$ be the image of $f(x)=x+\sin x$ about the line $x+y=0$. If the area bounded by $y=g(x)$, x-axis, $x=0$ and $x=2\pi$ is $A$, then $\dfrac{A}{\pi^2}$ is
\end{bankq}
\begin{bankq}[id=TOAS13-P2-MATH-S2-Q09, source={TOAS - Test 13 (24 Jul 2022)}, paper=2, qno=9, type=numerical, subtopic={Area between curves}, answer={3}, answer_src=solved, working={On $[1,\pi/2)$: $[x]=1$, $[\sin x]=0$, so $y=\sqrt{\sin1}$ and the area is $\sqrt{\sin1}\left(\dfrac\pi2-1\right)$: $a=2$, $b=1$, $a+b=3$}, difficulty=2, diff_parts={K1 S1 I0 C0 F1}, diff_note={concepts: greatest integer function in an area | steps: on [1, pi/2) [x] = 1 and [sin x] = 0 so y is constant; area times width}]
If area bounded by the curve $y=\sqrt{\sin[x]+[\sin x]}$, (where $[.]$ G.I.F.) and lines $x=1$ and $x=\dfrac\pi2$ and x-axis is equal to $\sqrt{\sin1}\left(\dfrac\pi a-b\right)$, then $a+b=$
\end{bankq}
\begin{bankq}[id=TOAS13-P2-MATH-S2-Q10, source={TOAS - Test 13 (24 Jul 2022)}, paper=2, qno=10, type=numerical, subtopic={Area between curves}, answer={4}, answer_src=solved, working={The two modulus conditions give the square with vertices $(\pm1,0),(0,\pm1)$ of area $2$; $2x^2+2y^2\ge1$ removes the inscribed circle of radius $1/\sqrt2$ (area $\pi/2$). Area $=2-\dfrac\pi2$: $a=2$, $b=2$, $a+b=4$}, difficulty=3, diff_parts={K1 S2 I0 C0 F1}, diff_note={concepts: area of a region between a square and a circle | steps: square from the modulus inequalities has area 2; circle of radius 1/sqrt 2 is inscribed; subtract}]
The area bounded by the curves $|x+y|\le1$, $|y-x|\le1$ and $2x^2+2y^2\ge1$ is $\left|a-\dfrac\pi b\right|$ then $a+b$ is equal to
\end{bankq}
\begin{bankq}[id=TOAS14-P1-MATH-S1-Q04, source={TOAS - Test 14 (07 Aug 2022)}, paper=1, qno=4, type=multi, subtopic={Area between curves}, answer={A,C}, answer_src=key, difficulty=2, diff_parts={K1 S1 I0 C0 F1}, diff_note={concepts: area under a parabola | steps: area is 4 a\^{}2 in magnitude; a can be negative since the area is a magnitude}]
Area enclosed by parabola $ay=3(a^2-x^2)$ and x-axis is 64, then value of `$a$' can be -
\begin{qoptions}
\opt{A}{4}
\opt{B}{2}
\opt{C}{$-4$}
\opt{D}{$-2$}
\end{qoptions}
\end{bankq}
\begin{bankq}[id=TOAS14-P1-MATH-S1-Q07, source={TOAS - Test 14 (07 Aug 2022)}, paper=1, qno=7, type=multi, subtopic={Area between curves}, answer={A,B,C}, answer_src=key, difficulty=5, diff_parts={K2 S2 I1 C0 F1}, diff_note={concepts: periodicity and symmetry of an area, integral bounds | steps: period pi/4 reduces S to 2 integral of e\^{}(sin t); the integrand is at least 1; sin t <= t gives the upper bound}]
Let $S$ be the area bounded by $y=e^{|\cos4x|}$, $x=0$, $y=0$ and $x=\pi$. Then
\begin{qoptions}
\opt{A}{$S=2\displaystyle\int_0^{\pi/2}e^{\sin t}\,dt$}
\opt{B}{$S>\dfrac\pi2$}
\opt{C}{$S<2\left(e^{\frac\pi2}-1\right)$}
\opt{D}{$S\le1$}
\end{qoptions}
\end{bankq}
\begin{bankq}[id=TOAS14-P1-MATH-S3-Q04, source={TOAS - Test 14 (07 Aug 2022)}, paper=1, qno=4, type=numerical, subtopic={Area between curves}, answer={5}, answer_src=key, difficulty=6, diff_parts={K2 S3 I0 C1 F1}, diff_note={concepts: minimum of an exponential curve, area between a curve and a chord | steps: minimum at x = ln 3; coordinates of P and Q; trapezoid under the chord; integral under the curve; subtract; match k}]
The curve $y=e^{2x}-18x+15$ crosses the y-axis at $P$ and has point of minima at $Q$. If the area bounded by the curve and chord $PQ$ is $(k\ln3-k+1)$ then $k$ is equal to (use $\ln3=1.1$)
\end{bankq}
\begin{bankq}[id=TOAS15-P1-MATH-S2-Q11, source={TOAS - Test 15 (22 Aug 2022)}, paper=1, qno=11, type=multi, subtopic={Area between curves}, answer={A,B}, answer_src=key, difficulty=3, diff_parts={K1 S1 I0 C1 F1}, diff_note={concepts: area between two curves | steps: intersection at x = +-1; integrate the difference}]
Area of region bounded between the curves $y=\dfrac1{1+x^2}$ and $y=1-\dfrac1{1+x^2}$ between its points of intersection is
\begin{qoptions}
\opt{A}{$\pi-2$}
\opt{B}{an irrational number}
\opt{C}{an integer}
\opt{D}{4}
\end{qoptions}
\end{bankq}
\begin{bankq}[id=TOAS15-P1-MATH-S3-Q05, source={TOAS - Test 15 (22 Aug 2022)}, paper=1, qno=5, type=numerical, subtopic={Area between curves}, answer={0.74}, answer_src=solved, working={$A=\int_0^1x(e^x-e^{-x})dx=\left[x(e^x+e^{-x})\right]_0^1-\int_0^1(e^x+e^{-x})dx=(e+e^{-1})-(e-e^{-1})=2/e\approx0.74$}, difficulty=2, diff_parts={K1 S1 I0 C0 F1}, diff_note={concepts: area between two curves | steps: difference x(e\^{}x - e\^{}-x); integrate by parts}]
The area bounded by the curves $y=xe^x$, $y=xe^{-x}$ and the lines $x=0$ and $x=1$ is (in \textit{square units}) :
\end{bankq}
\begin{bankq}[id=TOAS15-P2-MATH-S2-Q06, source={TOAS - Test 15 (22 Aug 2022)}, paper=2, qno=6, type=numerical, subtopic={Area between curves}, answer={4.00}, answer_src=key, difficulty=7, diff_parts={K2 S2 I2 C1 F1}, diff_note={concepts: area with an inverse function, lines | steps: reflect to relate the curve and its inverse; boundaries; integrate over the region; read k}]
If area of the region whose upper left boundary is enclosed by the curve $f(x)=\sin x+x+2$, lower left by the line $x+y=2$, lower right by $f^{-1}(x)$ and upper right by the line $x+y=2\pi+2$ is $k(\pi+1)$ then find the value of $k$
\end{bankq}
\begin{bankq}[id=TOAS16-P2-MATH-S4-Q07, source={TOAS - Test 16 (04 Sep 2022)}, paper=2, qno=7, type=single, subtopic={Area between curves}, answer={200}, answer_src=key, difficulty=2, diff_parts={K1 S1 I1 C0 F0}, diff_note={concepts: area under two curves | steps: split at the intersection; add the areas}]
If area bounded on the left by $y$-axis, below by $x$-axis, right by $x=\dfrac\pi2$ above left by $y=\cos x$ and above right by $y=\sin x$ is $K$ then value of $100K^2$ is -
\end{bankq}
\begin{bankq}[id=TOAS17-P2-MATH-S1-Q15, source={TOAS - Test 17 (18 Sep 2022)}, paper=2, qno=15, type=single, subtopic={Area between curves}, answer={C}, answer_src=key, difficulty=4, diff_parts={K1 S2 I1 C1 F0}, diff_note={concepts: area bounded by two parabolas and a line | steps: intersections; sketch the region; integrate; add}]
Area bounded by the curves $y=-x^2+6x-5$, $y=-x^2+4x-3$ and $y=3x-15$ is
\begin{qoptions}
\opt{A}{$\dfrac{73}2$}
\opt{B}{$\dfrac{73}3$}
\opt{C}{$\dfrac{73}6$}
\opt{D}{73}
\end{qoptions}
\end{bankq}
\begin{bankq}[id=TOAS18-P1-MATH-S3-Q02, source={TOAS - Test 18 (30 Oct 2022)}, paper=1, qno=2, type=numerical, subtopic={Area between curves}, answer={72}, answer_src=solved, working={Differentiating the relation gives $f\prime=\sqrt f$ so $f=x^2/4$; area between the parabola and $y=ax+9$ is smallest at $a=0$: $\tfrac1{24}(12)^3=72$}, difficulty=7, diff_parts={K2 S2 I2 C1 F1}, diff_note={concepts: functional equation, area between parabola and a line | steps: derive f' = sqrt f; f = x\^{}2/4; area formula; minimise over a}]
The minimum area bounded by the differentiable function $y=f(x)$ and $y=ax+9$ ($a\in R$) where $f$ satisfies the relation $f(x+y)=f(x)+f(y)+y\sqrt{f(x)}$, $x,y\in R$ and $f'(0)=0$ is
\end{bankq}
\begin{bankq}[id=TOAS18-P2-MATH-S2-Q10, source={TOAS - Test 18 (30 Oct 2022)}, paper=2, qno=10, type=multi, subtopic={Area between curves}, answer={A,C}, answer_src=key, difficulty=5, diff_parts={K1 S2 I1 C1 F1}, diff_note={concepts: area between two parabolas | steps: intersections; symmetry; integrate; divisibility}]
If area bounded by $y^2=8|x|$ and $x^2=3|y|$ is $\Delta$, then $\Delta$ is divisble is
\begin{qoptions}
\opt{A}{4}
\opt{B}{6}
\opt{C}{8}
\opt{D}{12}
\end{qoptions}
\end{bankq}

\section*{TAAS papers}
\begin{bankq}[id=TAAS02-P1-MATH-S3-Q11, source={TAAS - Test 02 (30 Jul 2023)}, paper=1, qno=11, type=single, subtopic={Area between curves}, answer={B}, answer_src=key, difficulty=5, diff_parts={K1 S3 I1 C1 F0}, diff_note={concepts: areas of various regions | steps: cycloid area; tangent triangle; hyperbola triangle; maximum area; match}]
Match the statement in List-I with given value in List-II
\begin{center}\footnotesize\setlength{\tabcolsep}{2pt}
\begin{tabular}{|L{0.6cm}L{4.4cm}|L{0.7cm}L{0.4cm}|}\hline
\multicolumn{2}{|c|}{List-I}&\multicolumn{2}{c|}{List-II}\\\hline
(P)&If area bounded by $x=(\theta-\sin\theta)$; $y=(1-\cos\theta)$; $0\leq\theta\leq2\pi$ is $\lambda\pi$; then the value of $\lambda$ can be equal to&(1)&6\\\hline
(Q)&If the least area of triangle formed by tangent to the circle $x^2+y^2=1$ and $x=0$, $y=0$ is $A$, then $A$ is coprime with&(2)&3\\\hline
(R)&Area of triangle formed by any tangent to the curve $y=\dfrac{2x^2-x-1}{2x^2+5x+2}$ $\left(x\neq-2,-\dfrac12\right)$ and pair of lines $xy-x+2y=2$ is&(3)&80\\\hline
(S)&If the maximum area bounded by the curves $x^2=4ay$, $y=ax$, $y=\dfrac xa$, $1\leq a\leq2$ is $\alpha$, then value of $(\alpha+a)$ is&(4)&84\\\hline
&&(5)&86\\\hline
\end{tabular}\end{center}
\begin{qoptions}
\opt{A}{P $\to$ 2;Q $\to$ 1,2;R $\to$ 1,4,5;S $\to$ 1,2}
\opt{B}{P $\to$ 2;Q $\to$ 1,2,3,4,5;R $\to$ 1;S $\to$ 4}
\opt{C}{P $\to$ 2;Q $\to$ 1,2,3,4,5;R $\to$ 2,5;S $\to$ 1,4}
\opt{D}{P $\to$ 3;Q $\to$ 2;R $\to$ 4;S $\to$ 4}
\end{qoptions}
\end{bankq}
\begin{bankq}[id=TAAS02-P2-MATH-S1-Q04, source={TAAS - Test 02 (30 Jul 2023)}, paper=2, qno=4, type=single, subtopic={Area between curves}, answer={B}, answer_src=key, difficulty=3, diff_parts={K1 S2 I0 C1 F0}, diff_note={concepts: areas under rational curves | steps: area under 1/(1+x\^{}2); intersection of the curves; area between}, also_in={TNAS - Test 06 (30 Jul 2023) P2 Q4}]
If $C_1\equiv y=\dfrac1{1+x^2}$ and $C_2\equiv y=\dfrac{x^2}2$ be two curve lying in $XY$ plane. Then
\begin{qoptions}
\opt{A}{area bounded by curve $y=\dfrac1{1+x^2}$ and $y=0$ is $2\pi$}
\opt{B}{area bounded by $C_1$ and $C_2$ is $\dfrac\pi2-\dfrac13$}
\opt{C}{area bounded by $C_1$ and $C_2$ is $1-\dfrac\pi2$}
\opt{D}{area bounded by curve $y=\dfrac1{1+x^2}$ and $x$-axis is $\dfrac\pi2$}
\end{qoptions}
\end{bankq}
\begin{bankq}[id=TAAS02-P2-MATH-S4-Q06, source={TAAS - Test 02 (30 Jul 2023)}, paper=2, qno=6, type=numerical, subtopic={Area between curves}, answer={3}, answer_src=key, difficulty=5, diff_parts={K1 S2 I1 C1 F1}, diff_note={concepts: areas cut by a parabola | steps: tangent at the origin; intersections; two areas; ratio}, also_in={TNAS - Test 06 (30 Jul 2023) P2 Q6}]
The parabola $y^2=x-y$ divides the triangle formed by the tangent to it at the origin and the straight line $x+y=0$ and $y=-1$ into two parts whose areas are $A_1$ and $A_2$ (Given $A_1<A_2$). If $\dfrac{A_1}{A_2}=\dfrac pq$ where $p,q\in N$ and $\dfrac pq$ is in simplest form then the value of $(p+q)$ is :
\end{bankq}

\section*{TMAS papers}
\begin{bankq}[id=TMAS04-P2-MATH-S2-Q06, source={TMAS - Test 04 (19 Nov 2023)}, paper=2, qno=6, type=single, subtopic={Area between curves}, answer={8}, answer_src=key, difficulty=5, diff_parts={K1 S2 I2 C1 F0}, diff_note={concepts: area between curves | steps: regions from the modulus condition; hyperbola; intersection; integral}]
The area of the region which consists of all the points satisfying the conditions $|x-y|+|x+y|\le8$, $xy\ge2$, $x\ge0$ and $y\ge0$ is equal to - (take $\ln2=0.69$)
\end{bankq}

\section*{TNAS papers}
\begin{bankq}[id=TNAS06-P1-MATH-S3-Q11, source={TNAS - Test 06 (30 Jul 2023)}, paper=1, qno=11, type=single, subtopic={Area between curves}, answer={B}, answer_src=key, difficulty=6, diff_parts={K1 S3 I2 C2 F0 adj-1}, diff_note={concepts: areas in parametric and tangent problems | steps: area under a cycloid; least area of a tangent triangle; tangent to the rational function; maximum area; match values; check; select}]
Match the statement in List-I with given value in List-II\par
{\footnotesize\setlength{\tabcolsep}{2pt}\begin{tabular}{|c|p{0.62\linewidth}|c|p{0.06\linewidth}|}
\hline
\multicolumn{2}{|c|}{\textbf{List-I}} & \multicolumn{2}{c|}{\textbf{List-II}}\\ \hline
(P) & If area bounded by $x=(\theta-\sin\theta)$; $y=(1-\cos\theta)$; $0\le\theta\le2\pi$ is $\lambda\pi$; then the value of $\lambda$ can be equal to & (1) & 6\\ \hline
(Q) & If the least area of triangle formed by tangent to the circle $x^2+y^2=1$ and $x=0$, $y=0$ is $A$, then $A$ is coprime with & (2) & 3\\ \hline
(R) & Area of triangle formed by any tangent to the curve $y=\dfrac{2x^2-x-1}{2x^2+5x+2}\left(x\ne-2,-\dfrac12\right)$ and pair of lines $xy-x+2y=2$ is & (3) & 80\\ \hline
(S) & If the maximum area bounded by the curves $x^2=4ay$, $y=ax$, $y=\dfrac xa$, $1\le a\le2$ is $\alpha$, then value of $(\alpha+a)$ is & (4) & 84\\ \hline
 & & (5) & 86\\ \hline
\end{tabular}}\par
\begin{qoptions}
\opt{A}{$P\to2;Q\to1,2;R\to1,4,5;S\to1,2$}
\opt{B}{$P\to2;Q\to1,2,3,4,5;R\to1;S\to4$}
\opt{C}{$P\to2;Q\to1,2,3,4,5;R\to2,5;S\to1,4$}
\opt{D}{$P\to3;Q\to2;R\to4;S\to4$}
\end{qoptions}
\end{bankq}
\begin{bankq}[id=TNAS07-P1-MATH-S1-Q15, source={TNAS - Test 07 (20 Aug 2023)}, paper=1, qno=15, type=single, subtopic={Area between curves}, answer={C}, answer_src=key, difficulty=7, diff_parts={K1 S3 I2 C2 F0}, diff_note={concepts: area between two parabolic curves as a function of a parameter | steps: intersection points; integrate the difference; simplify; maximise; check statements}]
Consider the two curves $y=x-\lambda x^2$ and $y=\dfrac{x^2}\lambda$, ($\lambda>0$).\par
\textbf{Statement-1 :} The area bounded between the curves is maximum when $\lambda=1$.\par
and\par
\textbf{Statement-2 :} The area bounded between the curves is $\dfrac{\lambda^2}{(1+\lambda^2)^2}$ square units.
\begin{qoptions}
\opt{A}{Statement-1 is True, Statement-2 is True.}
\opt{B}{Statement-1 is False, Statement-2 is False.}
\opt{C}{Statement-1 is True, Statement-2 is False.}
\opt{D}{Statement-1 is False, Statement-2 is True.}
\end{qoptions}
\end{bankq}
\begin{bankq}[id=TNAS07-P2-MATH-S2-Q10, source={TNAS - Test 07 (20 Aug 2023)}, paper=2, qno=10, type=single, subtopic={Area between curves}, answer={D}, answer_src=key, difficulty=8, diff_parts={K2 S3 I2 C2 F0}, diff_note={concepts: areas, integral and slope matched | steps: area between the trigonometric curves; parabola area; sum of cosines integral; slope of the line; match}]
{\small\setlength{\tabcolsep}{2pt}\begin{tabular}{|c|p{0.58\linewidth}|c|c|}
\hline
\multicolumn{2}{|c|}{\textbf{List-I}} & \multicolumn{2}{c|}{\textbf{List-II}}\\ \hline
(I) & The area bounded by the curves $y=\tan^2x$, $y=\sec^2x$, $x=-\dfrac\pi4$ and $x=\dfrac\pi4$ is given by & (P) & $\dfrac{16}3$\\ \hline
(II) & The area bounded by $y=x^2-4$ with the co-ordinate axes is & (Q) & $\dfrac\pi2$\\ \hline
(III) & $\displaystyle\int_0^{\pi/2}\dfrac{\sin6x}{\sin x}dx$ is equal to & (R) & $\dfrac{32}3$\\ \hline
(IV) & The slope of line which cuts-off equal intercept on co-ordinate axes can be & (S) & $\dfrac{26}{15}$\\ \hline
 & & (T) & $1$\\ \hline
\end{tabular}}\par
\begin{qoptions}
\opt{A}{$\mathrm I\to P,Q,R,S,T;\allowbreak \mathrm{II}\to S;\allowbreak \mathrm{III}\to P,R;\allowbreak \mathrm{IV}\to Q$}
\opt{B}{$\mathrm I\to Q;\allowbreak \mathrm{II}\to P,R;\allowbreak \mathrm{III}\to P,Q,R,S,T;\allowbreak \mathrm{IV}\to S$}
\opt{C}{$\mathrm I\to P,R;\allowbreak \mathrm{II}\to Q;\allowbreak \mathrm{III}\to S;\allowbreak \mathrm{IV}\to P,Q,R,S,T$}
\opt{D}{$\mathrm I\to Q;\allowbreak \mathrm{II}\to P,R;\allowbreak \mathrm{III}\to S;\allowbreak \mathrm{IV}\to P,Q,R,S,T$}
\end{qoptions}
\end{bankq}
\end{document}
